\documentclass[11pt]{amsart}
\usepackage[a4paper,margin=1in]{geometry}
\usepackage{amsmath,amssymb,mathtools}
\usepackage{enumitem}
\usepackage{microtype}
\usepackage[hidelinks]{hyperref}

\newtheorem{theorem}{Theorem}[section]
\newtheorem{proposition}[theorem]{Proposition}
\newtheorem{lemma}[theorem]{Lemma}
\newtheorem{corollary}[theorem]{Corollary}

\usepackage[nameinlink,noabbrev]{cleveref}

\crefname{theorem}{theorem}{theorems}
\Crefname{theorem}{Theorem}{Theorems}
\crefname{proposition}{proposition}{propositions}
\Crefname{proposition}{Proposition}{Propositions}
\crefname{lemma}{lemma}{lemmas}
\Crefname{lemma}{Lemma}{Lemmas}
\crefname{corollary}{corollary}{corollaries}
\Crefname{corollary}{Corollary}{Corollaries}

\newcommand{\Z}{\mathbb Z}
\newcommand{\Q}{\mathbb Q}
\newcommand{\R}{\mathbb R}
\newcommand{\C}{\mathbb C}
\newcommand{\PP}{\mathbb P}
\newcommand{\HH}{\mathfrak H}
\newcommand{\cO}{\mathcal O}
\newcommand{\cV}{\mathcal V}
\newcommand{\cH}{\mathcal H}
\newcommand{\GL}{\operatorname{GL}}
\newcommand{\SL}{\operatorname{SL}}
\newcommand{\im}{\operatorname{im}}
\newcommand{\rk}{\operatorname{rk}}
\newcommand{\Hom}{\operatorname{Hom}}
\newcommand{\Spec}{\operatorname{Spec}}
\newcommand{\vol}{\operatorname{vol}}
\newcommand{\wt}{\widetilde}
\newcommand{\wh}{\widehat}
\newcommand{\ol}{\overline}
\newcommand{\eps}{\varepsilon}

\title[A compact complex threefold diffeomorphic to $S^6$]
{A compact complex threefold diffeomorphic to the six-sphere}
\author[ChatGPT Sol 5.6 PRO]{Streamlining by ChatGPT Sol 5.6 PRO of Anthropic's original proof, as communicated by Levent Alpoge}
\date{}

\begin{document}
\begin{abstract}
We construct a compact complex threefold fibred by complex two-tori over
$\mathbb P^1$.  The fibration has one reduced, nonnormal normal crossings fibre and
two multiple fibres of multiplicities $3$ and $4$.  We prove that
the total space is simply connected and compute its integral cohomology by
the Leray spectral sequence.  The local calculation at the nonnormal fibre
is carried out integrally by nearby cycles.  The total space is an integral
homology six-sphere and hence is diffeomorphic to the standard $S^6$.
\end{abstract}

\maketitle

\section{Statement and conventions}

\begin{theorem}\label{thm:main}
There is a compact connected complex manifold $X$ of complex dimension
three whose underlying smooth manifold is diffeomorphic to the standard
six-sphere.  Consequently $S^6$ admits an integrable complex structure.
\end{theorem}

The proof presented here is a streamlining by ChatGPT Sol 5.6 PRO of an
original proof produced by an AI model internal to Anthropic, as communicated
by Levent Alpoge and circulated as \cite{AnthropicProof}; all arguments used
for the theorem are included here.

All lattices are written additively.  Exterior products are abbreviated by
juxtaposition, for example $\gamma u=\gamma\wedge u$.  Paths are composed
from left to right.  All meridians are oriented clockwise; based meridians
$a_1,a_2,a_0$ about $p_1,p_2,p_0$, respectively, are chosen so that
$a_1a_2a_0=1$.


\section{Integral monodromy and the base orbifold}

\subsection{The lattice representation}
Let $V=\Z^4$ with ordered basis $(\gamma,u,w,\delta)$ and let
$\Lambda=V^*$ with dual basis
$(\wh\gamma,\wh u,\wh w,\wh\delta)$.  Evaluation on $\gamma$ is denoted
by the same letter:
\[
 \gamma(a\wh\gamma+b\wh u+c\wh w+d\wh\delta)=a.
\]
Set
\[
T_1=
\begin{pmatrix}
1&0&-6&2\\
0&-1&1&1\\
0&-1&0&1\\
0&0&0&1
\end{pmatrix},
\qquad
T_2=
\begin{pmatrix}
1&6&0&-3\\
0&0&-1&1\\
0&1&0&0\\
0&0&0&1
\end{pmatrix}.
\]
For $T\in\GL(V)$ write $A(T)=(T^{-1})^t\in\GL(\Lambda)$, and put
$A_j=A(T_j)$.

\begin{proposition}[Integral linear algebra]\label{prop:linear-algebra}
The following statements hold.
\begin{enumerate}[label=\textup{(\roman*)}]
\item $T_1$ and $T_2$ have exact orders $3$ and $4$, respectively, and
\[
T_0:=(T_1T_2)^{-1}
 =\begin{pmatrix}
1&0&0&1\\
0&1&-1&0\\
0&0&1&0\\
0&0&0&1
\end{pmatrix}=I+N,
\qquad N^2=0,
\]
with
\[
N\gamma=Nu=0,\qquad Nw=-u,\qquad N\delta=\gamma.
\]
\item In the dual basis,
\[
\begin{aligned}
A_1\wh\gamma&=\wh\gamma+6\wh u-6\wh w-2\wh\delta,
&A_1\wh u&=-\wh w+\wh\delta,
&A_1\wh w&=\wh u-\wh w,
&A_1\wh\delta&=\wh\delta,\\
A_2\wh\gamma&=\wh\gamma-6\wh w+3\wh\delta,
&A_2\wh u&=\wh w,
&A_2\wh w&=-\wh u+\wh\delta,
&A_2\wh\delta&=\wh\delta.
\end{aligned}
\]
Writing $M_0=A(T_0)$, one has
\[
M_0\wh\gamma=\wh\gamma-\wh\delta,
\quad M_0\wh u=\wh u+\wh w,
\quad M_0\wh w=\wh w,
\quad M_0\wh\delta=\wh\delta,
\]
and $A_1A_2M_0=I$.
\item Put
\[
 \Lambda_{\mathrm{tor}}=\langle\wh w,\wh\delta\rangle,
 \qquad \ol\Lambda=\Lambda/\Lambda_{\mathrm{tor}}.
\]
Then
\[
\ker(M_0-I)=\im(M_0-I)=\Lambda_{\mathrm{tor}},
\]
and the induced map
\[
 B_0:\ol\Lambda\longrightarrow\Lambda_{\mathrm{tor}},
 \qquad B_0\ol{\wh\gamma}=-\wh\delta,
 \quad B_0\ol{\wh u}=\wh w,
\]
is unimodular.
\item The vectors
\[
 \eps=\wh\gamma+2\wh u-4\wh w,
 \qquad
 \eps'=\wh\gamma+3\wh u-3\wh w
\]
satisfy
\[
 \Lambda^{A_1}=\langle\eps,\wh\delta\rangle,
 \qquad
 \Lambda^{A_2}=\langle\eps',\wh\delta\rangle.
\]
We fix the twist vectors
\[
 v_1=\eps,
 \qquad v_2=-\eps'.
\]
Thus $\gamma(v_1)=1$ and $\gamma(v_2)=-1$.
\item If $G=\langle T_1,T_2\rangle$, then
\[
 V^G=\Z\gamma,
 \qquad
 \sum_{g\in G}\im(g-I)=\langle\gamma,u,w\rangle,
\]
while
\[
 \Lambda^G=\Z\wh\delta,
 \qquad
 \sum_{g\in G}\im(A(g)-I)=\ker\gamma.
\]
In particular the coinvariants $V_G$ and $\Lambda_G$ are free of rank one.
The smallest $\langle A_1,A_2\rangle$-invariant subgroup of $\Lambda$
containing $\Lambda_{\mathrm{tor}}$ is $\ker\gamma$.
\item Set
\[
 q=uw+6\gamma\delta,
 \qquad
 \vol=\gamma uw\delta.
\]
Then
\[
 (\textstyle\bigwedge^kV)^G=
 \begin{cases}
 \Z,&k=0,\\
 \Z\gamma,&k=1,\\
 \Z q,&k=2,\\
 \Z\gamma uw,&k=3,\\
 \Z\vol,&k=4,
 \end{cases}
\]
and the coinvariants are
\[
 V_G=\Z[\delta],
 \quad
 (\textstyle\bigwedge^2V)_G=\Z[\gamma\delta],
 \quad [uw]=6[\gamma\delta],
 \quad
 (\textstyle\bigwedge^3V)_G=\Z[uw\delta],
 \quad
 (\textstyle\bigwedge^4V)_G=\Z[\vol].
\]
Finally
\[
\begin{aligned}
 V^{T_0}&=\langle\gamma,u\rangle,\\
 (\textstyle\bigwedge^2V)^{T_0}
 &=\langle\gamma u,\gamma\delta,uw,\gamma w-u\delta\rangle,\\
 (\textstyle\bigwedge^3V)^{T_0}
 &=\langle\gamma uw,\gamma u\delta\rangle.
\end{aligned}
\]
All displayed invariant subgroups are saturated, and all displayed
coinvariant groups are torsion-free.
\end{enumerate}
\end{proposition}

\begin{proof}
Direct multiplication gives $T_1^3=T_2^4=I$ and the displayed formula
for $T_0$.  The formulas for $A_1,A_2,M_0$ follow by taking inverse
transposes.  For
$\lambda=a\wh\gamma+b\wh u+c\wh w+d\wh\delta$,
\[
 (M_0-I)\lambda=-a\wh\delta+b\wh w,
\]
which proves (iii).  Solving $(A_j-I)\lambda=0$ gives
$b=2a,c=-4a$ for $j=1$, and $b=3a,c=-3a$ for $j=2$, proving (iv).

For (v), solving $T_jx=x$ gives
\[
 V^{T_1}=\langle\gamma,2u+w+3\delta\rangle,
 \qquad
 V^{T_2}=\langle\gamma,u+w+2\delta\rangle,
\]
whose intersection is $\Z\gamma$.  The images of $T_j-I$ contain
$u,w,2\gamma,3\gamma$, hence generate
$\langle\gamma,u,w\rangle$.  Dually, the images of $A_j-I$ generate
$\langle\wh u,\wh w,\wh\delta\rangle=\ker\gamma$.  The last assertion
follows because an invariant subgroup containing $\wh w$ contains
$A_1\wh w=\wh u-\wh w$, hence $\wh u$ and $\wh\delta$.

The exterior-power calculations in (vi) are given in
\cref{sec:matrix-tables}.  The Smith normal forms there have only unit
nonzero elementary divisors, which proves the asserted saturation.  The
$T_0$-invariants also follow directly from
$Nw=-u$, $N\delta=\gamma$, and $N\gamma=Nu=0$.
\end{proof}

\subsection{The orbifold curve}
Let $B=\PP^1$ with affine coordinate $t$, and put
\[
 p_1=0,
 \qquad p_2=1,
 \qquad p_0=\infty,
 \qquad t_c=t^{-1},
 \qquad B^\circ=B\setminus\{p_0,p_1,p_2\}.
\]
Let $\mathcal B$ be the orbifold with underlying curve $B$, cone orders
$3$ and $4$ at $p_1,p_2$, and a cusp at $p_0$.

\begin{proposition}[Uniformization]\label{prop:uniformization}
There are a copy $\HH_z$ of the upper half-plane, a Fuchsian group
$\Delta<\mathrm{PSL}_2(\R)$, and a holomorphic quotient map
\[
 \pi:\HH_z\longrightarrow B\setminus\{p_0\}
\]
with the following properties.
\begin{enumerate}[label=\textup{(\roman*)}]
\item $\HH_z/\Delta\cong B\setminus\{p_0\}$ as orbifolds.
\item There are $z_j\in\pi^{-1}(p_j)$ and generators $g_1,g_2$ such that
\[
 \Delta=\langle g_1,g_2\mid g_1^3=g_2^4=1\rangle,
 \qquad g_0=(g_1g_2)^{-1}
\]
is primitive parabolic, and
\[
 g_j'(z_j)=e^{-2\pi i/m_j},
 \qquad (m_1,m_2)=(3,4).
\]
\item For sufficiently small $r>0$, the component
$U_r\subset\pi^{-1}(0<|t_c|<r)$ accumulating at the fixed point of
$g_0$ is simply connected, invariant under $\langle g_0\rangle$, disjoint
from all its other $\Delta$-translates, and
\[
 U_r/\langle g_0\rangle\cong\{0<|t_c|<r\}.
\]
\end{enumerate}
\end{proposition}

\begin{proof}
This is the standard uniformization of the hyperbolic triangle orbifold of
signature $(0;3,4,\infty)$; see \cite[Chs.~7,10]{Beardon}.  Choose a hyperbolic triangle with angles
$\pi/3,\pi/4,0$ and take the orientation-preserving subgroup of its
reflection group.  If the vertices are ordered counterclockwise, the
inverses of the positive rotations about the finite vertices give
$g_1,g_2$ with the stated derivatives, while $(g_1g_2)^{-1}$ is the
primitive cusp generator.  The standard disjoint-horodisc theorem for a
cofinite Fuchsian group gives (iii).
\end{proof}

The homomorphism $\rho_V:\Delta\to\SL(V)$ is defined by
$\rho_V(g_j)=T_j$.  The monodromy on $\Lambda$ will be the contragredient
representation $g_j\mapsto A_j$, $g_0\mapsto M_0$.


\section{The period family over the punctured base}

Let $\rho=e^{\pi i/3}$.  We use the classical modular function
$j:\HH\to\C$, normalized by $j(i)=1728$, and the modular forms
$E_4,E_6,\Delta_{\mathrm{mod}}$ of weights $4,6,12$ \cite[Ch.~VII]{Serre}.

\subsection{The three period functions}

\begin{theorem}[Periods]\label{thm:periods}
There are holomorphic functions
\[
 \tau:\HH_z\to\HH,
 \qquad
 \mu,\beta:\HH_z\to\C
\]
satisfying
\begin{align}
 j(\tau)&=1728\,t\circ\pi,\label{eq:tau-j}\\
 \tau\circ g_1&=\frac{\tau-1}{\tau},
 &\tau\circ g_2&=-\frac1\tau,
 &\tau(z_1)&=\rho,
 &\tau(z_2)&=i,\label{eq:tau-law}\\
 \mu\circ g_1&=\frac{1-\mu}{\tau},
 &\mu\circ g_2&=1+\frac\mu\tau,
 \label{eq:mu-law}\\
 \beta\circ g_1&=\beta+2-\frac{6(1-\mu)^2}{\tau},
 &\beta\circ g_2&=\beta-3-\frac{6\mu^2}{\tau}.
 \label{eq:beta-law}
\end{align}
These laws imply
\begin{equation}\label{eq:cusp-period-laws}
 \tau\circ g_0=\tau-1,
 \qquad
 \mu\circ g_0=\mu,
 \qquad
 \beta\circ g_0=\beta+1.
\end{equation}
On the distinguished cusp component $U_r$, the functions $\mu$ and
$\beta+\tau$ are bounded and descend holomorphically across $t_c=0$.
The functions $\tau$ and $\mu$ are uniquely determined by these conditions,
whereas the set of functions $\beta$ satisfying the stated conditions is
an affine space under the additive constants.
\end{theorem}

\begin{proof}
\smallskip
\noindent\emph{The function $\tau$.}
Let $\mathcal M$ denote the modular orbifold, whose coarse coordinate is
$j$ and whose orbifold orders at $0$, $1728$, and $\infty$ are
$3$, $2$, and $\infty$; the map $j:\HH\to\mathcal M$ is its universal
orbifold covering.  The coarse map
\[
 \varphi:B\longrightarrow\PP^1,
 \qquad \varphi(t)=1728t,
\]
is a holomorphic orbifold morphism from
$(B;3,4,\infty)$ to $\mathcal M$.  We verify the local lifts and their
isotropy homomorphisms.  Choose source uniformizers $s_1,s_2$ at
$p_1,p_2$ so that, after multiplying them by invariant holomorphic units,
\[
 t=s_1^3,\qquad t-1=s_2^4.
\]
Choose target uniformizers $\xi_1,\xi_2$ at the orbifold points $0,1728$ so
that $j=\xi_1^3$ and $j-1728=\xi_2^2$, again after multiplication by
isotropy-invariant units.
The map $\varphi$ then has local lifts of the form
\[
 \xi_1=s_1\cdot(\text{unit}),
 \qquad
 \xi_2=s_2^2\cdot(\text{unit}).
\]
Thus the order-three generator maps to
$S_1=\begin{psmallmatrix}1&-1\\1&0\end{psmallmatrix}$, whose derivative
at $\rho=e^{\pi i/3}$ is $e^{-2\pi i/3}$, while the order-four generator
maps to
$S=\begin{psmallmatrix}0&-1\\1&0\end{psmallmatrix}$: the action
$s_2\mapsto-is_2$ induces $\xi_2\mapsto-\xi_2$.  At the cusp, write $q=e^{2\pi i\zeta}$ for the standard modular
coordinate, where $\zeta\in\HH$.  Since
$j=q^{-1}(1+O(q))$, the equation $j=1728/t_c$ determines a local lift
\begin{equation}\label{eq:q-cusp}
 q=t_cu_1(t_c)
\end{equation}
with $u_1$ a holomorphic unit.  A clockwise loop about $t_c=0$ therefore
maps to the clockwise modular cusp generator
$T^{-1}=\begin{psmallmatrix}1&-1\\0&1\end{psmallmatrix}$.  Equivalently,
$g_1,g_2,g_0$ map to $S_1,S,T^{-1}$; indeed
$(S_1S)^{-1}=T^{-1}$.

The orbifold lifting theorem \cite[Chs.~7,10]{Beardon} now gives a holomorphic lift
$\tau:\HH_z\to\HH$ equivariant for the homomorphism
$g_1\mapsto S_1$, $g_2\mapsto S$.  It satisfies $\tau(z_1)=\rho$.
This equivariant lift is unique: two such lifts differ by an element of
$\mathrm{PSL}_2(\Z)$ centralizing both $S_1$ and $S$, and that common
centralizer is trivial.  The preceding isotropy homomorphisms give
\eqref{eq:tau-j}--\eqref{eq:tau-law}.  The local lift at $p_1$ has degree
one, so $\tau-\rho$ has order one at $z_1$; the local lift at $p_2$ has
degree two, so $\tau-i$ has order two at $z_2$.  On $U_r$, the cusp
coordinate in \eqref{eq:q-cusp} is $q=e^{2\pi i\tau}$.

Choose a branch of $h=(2\pi i)^{-1}\log u_1$ in
\eqref{eq:q-cusp} and set
\begin{equation}\label{eq:s-cusp}
 s=\tau-h(t_c\circ\pi).
\end{equation}
Then $e^{2\pi is}=t_c\circ\pi$ and $s\circ g_0=s-1$.  Both
$t_c\circ\pi:U_r\to\{0<|t_c|<r\}$ and the exponential map from a
half-plane are universal coverings with the same deck generator.  Hence
$s$ is a biholomorphism from $U_r$ onto a half-plane.

\smallskip
\noindent\emph{The function $\mu$.}
Define a multiplicative cocycle $\jmath$ on $\Delta\times\HH_z$ by
\[
 \jmath(g_1,z)=-\tau(z),
 \qquad
 \jmath(g_2,z)=\tau(z),
\]
and the cocycle rule.  The identities for the $g_1$- and $g_2$-orbits of
$\tau$ show that this is well defined and that $\jmath(g_0,z)=1$.
For an open set $U\subset B$, let $\mathcal L(U)$ be the holomorphic
functions $\nu$ on $\pi^{-1}(U\setminus\{p_0\})$ satisfying
\[
 \nu\circ g_1=-\frac\nu\tau,
 \qquad
 \nu\circ g_2=\frac\nu\tau,
\]
with the following additional condition when $p_0\in U$: on the
distinguished cusp component, $\nu$ descends as a function of $t_c$ and
extends holomorphically across $t_c=0$.  These spaces form a sheaf of
$\cO_B$-modules, denoted $\mathcal L$.
We claim
\begin{equation}\label{eq:L-minus-one}
 \mathcal L\cong\cO_B(-p_0)\cong\cO_{\PP^1}(-1).
\end{equation}
Indeed, since $\HH_z$ is simply connected and $E_6\circ\tau$ has zeros
of even order, choose a holomorphic square root and set
\[
 F=\left(\frac{E_4^2E_6^{1/2}}{\Delta_{\mathrm{mod}}}\right)\circ\tau.
\]
Choose lifts to $\mathrm{SL}_2(\Z)$ of the images of $g_1$ and
$g_2$ in $\mathrm{PSL}_2(\Z)$.  The numerator and denominator in the
definition of $F$ have total modular weights $11$ and $12$; hence their
quotient has weight $-1$.  The sign of the chosen square root of $E_6$
and the sign of an $\mathrm{SL}_2(\Z)$-lift combine into constants
$\chi_j\in\{\pm1\}$ for which
\[
 F\circ g_j=\frac{\chi_jF}{\jmath(g_j)}.
\]
Iterating around $g_1^3=1$ gives $\chi_1^3=1$, hence $\chi_1=1$.
At $z_2$, the function $E_6^{1/2}\circ\tau$ has a simple zero in the
linearizing coordinate $s_2$, and $s_2\circ g_2=-is_2$.  The weight-three
automorphy factor at $z_2$ is $\tau(z_2)^3=i^3=-i$, so comparison of the
linear terms gives $\chi_2=1$.  Thus $F$ satisfies the homogeneous laws
defining $\mathcal L$.

The function $F$ has order $2$ at $z_1$, order $1$ at $z_2$, and is
$t_c^{-1}$ times a unit at the cusp.  If $\nu$ is another local
homogeneous solution, then $\nu/F$ is $\Delta$-invariant and hence a
meromorphic function of $t$.  At $p_1$ a pole of $\nu/F$ would lift to a
pole of order at most $2$ in the coordinate $s_1$, hence of order at most
$2/3$ in $t$; no integral pole order is possible.  The same argument at
$p_2$ gives the bound $1/4$.  At $p_0$, regularity of $\nu$ and the simple
pole of $F$ force the quotient to vanish.  Thus
$\mathcal L_{p_0}=\mathfrak m_{p_0}F$ and $F$ generates elsewhere,
proving \eqref{eq:L-minus-one}.

The affine transformation laws \eqref{eq:mu-law} are consistent with
$g_1^3=g_2^4=1$.  Explicitly, iteration gives
\[
\begin{array}{c|ccc}
 &z&g_1z&g_1^2z\\ \hline
\tau&\tau&(\tau-1)/\tau&-1/(\tau-1)\\
\mu&\mu&(1-\mu)/\tau&(\tau-1+\mu)/(\tau-1)
\end{array}
\]
and
\[
\begin{array}{c|cccc}
 &z&g_2z&g_2^2z&g_2^3z\\ \hline
\tau&\tau&-1/\tau&\tau&-1/\tau\\
\mu&\mu&1+\mu/\tau&1-\tau-\mu&(1-\mu)/\tau.
\end{array}
\]
A direct substitution gives
\[
 \mu\circ(g_1g_2)=\mu,
 \qquad\text{hence}\qquad
 \mu\circ g_0=\mu.
\]
Consequently the local solutions form a torsor under $\mathcal L$.
On a disc avoiding the three special points, prescribe a holomorphic
function on one lift and extend it by the affine laws.  At the cusp the
identity $\mu\circ g_0=\mu$ makes the choice $\mu=0$ consistent on the
distinguished component.  On invariant discs at the elliptic points use
\[
 \mu_1=\frac{2-\tau}{3},
 \qquad
 \mu_2=\frac{1-\tau}{2}.
\]
Since $H^1(B,\mathcal L)=0$, the torsor has a global section; since
$H^0(B,\mathcal L)=0$, it is unique.  At the cusp it is bounded and hence
extends holomorphically.  Evaluating at the elliptic fixed points gives
\begin{equation}\label{eq:mu-values}
 \mu(z_1)=\frac1{1+\rho},
 \qquad
 \mu(z_2)=\frac{1-i}{2},
\end{equation}
and therefore
\[
 6(1-\mu(z_1))^2=2\tau(z_1),
 \qquad
 6\mu(z_2)^2=-3\tau(z_2).
\]

\smallskip
\noindent\emph{The function $\beta$.}
Put
\[
 \phi_1=2-\frac{6(1-\mu)^2}{\tau},
 \qquad
 \phi_2=-3-\frac{6\mu^2}{\tau}.
\]
Substitution of the preceding orbit tables gives
\begin{align*}
 \phi_1\circ g_1&=2-\frac{6(\tau-1+\mu)^2}{\tau(\tau-1)},
 &\phi_1\circ g_1^2&=2+\frac{6\mu^2}{\tau-1},\\
 \phi_2\circ g_2&=-3+\frac{6(\tau+\mu)^2}{\tau},
 &\phi_2\circ g_2^2&=-3-\frac{6(1-\tau-\mu)^2}{\tau},\\
 \phi_2\circ g_2^3&=-3+\frac{6(1-\mu)^2}{\tau}.
\end{align*}
A common-denominator calculation yields
\[
 \sum_{k=0}^2\phi_1\circ g_1^k=0,
 \qquad
 \sum_{k=0}^3\phi_2\circ g_2^k=0.
\]
Moreover \eqref{eq:mu-values} gives $\phi_j(z_j)=0$.  Substituting the
transformation laws once more gives
\[
 \beta\circ(g_1g_2)=\beta-1,
 \qquad
 \beta\circ g_0=\beta+1.
\]
Since $\tau\circ g_0=\tau-1$, the function $\beta+\tau$ is
$g_0$-invariant.  Hence the local solutions of \eqref{eq:beta-law}, with
$\beta+\tau$ extending at $p_0$, form a torsor under $\cO_B$.  At the
cusp the choice $\beta=-\tau$ is consistent with the induced $g_0$-law.
On a simply covered disc choose an arbitrary solution, and at $p_j$ use
\[
 \beta_j=\frac1{m_j}\sum_{k=0}^{m_j-1}k\,\phi_j\circ g_j^k.
\]
Since $H^1(B,\cO_B)=0$, a global solution exists, and all solutions differ
by constants.
\end{proof}

\subsection{Equivariance and the torus family}
Define
\[
 \Pi(z)=
 \begin{pmatrix}
 6\mu(z)&\tau(z)&1&0\\
 \beta(z)&\mu(z)&0&1
 \end{pmatrix}
 =[Z(z)\mid I_2].
\]
Regard its rows as the linear forms
\[
 \sigma_1=6\mu\gamma+\tau u+w,
 \qquad
 \sigma_2=\beta\gamma+\mu u+\delta
 \quad\text{in }V\otimes\C.
\]

\begin{proposition}[Equivariance]\label{prop:period-equivariance}
For the generators of $\Delta$ one has
\[
 R_1(z)=
 \begin{pmatrix}
 -1/\tau&0\\
 (1-\mu)/\tau&1
 \end{pmatrix},
 \qquad
 R_2(z)=
 \begin{pmatrix}
 1/\tau&0\\
 -\mu/\tau&1
 \end{pmatrix},
 \qquad R_0=I_2,
\]
and
\begin{equation}\label{eq:period-equivariance}
 \Pi(gz)=R_g(z)\Pi(z)M_g,
 \qquad M_g=\rho_V(g)^t.
\end{equation}
The matrices satisfy $R_{gh}(z)=R_g(hz)R_h(z)$.
\end{proposition}

\begin{proof}
Using the matrices of \cref{prop:linear-algebra},
\[
\begin{aligned}
T_1\sigma_1&=(6\mu-6)\gamma+(1-\tau)u-\tau w,\\
T_1\sigma_2&=(\beta+2)\gamma+(1-\mu)u+(1-\mu)w+\delta,\\
T_2\sigma_1&=(6\mu+6\tau)\gamma-u+\tau w,\\
T_2\sigma_2&=(\beta+6\mu-3)\gamma+u+\mu w+\delta.
\end{aligned}
\]
Substituting \eqref{eq:tau-law}--\eqref{eq:beta-law} gives
\[
 \sigma_1(g_1z)=-\tau^{-1}T_1\sigma_1,
 \quad
 \sigma_2(g_1z)=T_1\sigma_2+(1-\mu)\tau^{-1}T_1\sigma_1,
\]
and the analogous formulas for $g_2$.  For $g_0$ one obtains
$\sigma_i(g_0z)=T_0\sigma_i(z)$.  Evaluating on $\Lambda$ gives
\eqref{eq:period-equivariance}; the cocycle rule follows by uniqueness.
\end{proof}

\begin{corollary}[Choice of the additive constant]
\label{cor:nondegenerate-periods}
The function $\beta$ may be chosen so that
\begin{equation}\label{eq:D-negative}
 D:=\Im\beta-\frac{6(\Im\mu)^2}{\Im\tau}<0
 \qquad\text{on }\HH_z.
\end{equation}
For this choice, $L(z)=\Pi(z)\Lambda$ is a lattice in $\C^2$ for every
$z\in\HH_z$.
\end{corollary}

\begin{proof}
As a real linear map $\Lambda\otimes\R\to\C^2$,
\begin{equation}\label{eq:real-det}
 \det_{\R}\Pi(z)=-\det\Im Z(z)=\Im\tau(z)\,D(z).
\end{equation}
For both generators,
\[
 |\det_{\C}R_g(z)|^2=|\tau(z)|^{-2},
 \qquad
 \Im\tau(gz)=\Im\tau(z)|\tau(z)|^{-2}.
\]
Taking real determinants in \eqref{eq:period-equivariance} and using
$\det M_g=1$, comparison with \eqref{eq:real-det} therefore gives
$D(gz)=D(z)$.  Thus $D$ is $\Delta$-invariant.  Since $\tau,\mu,\beta$ are holomorphic at the
elliptic fixed points, it descends continuously to
$B\setminus\{p_0\}$.  At the cusp, $\beta+\tau$ and $\mu$ are bounded
while $\Im\tau\to\infty$, so $D\to-\infty$.  Hence $D$ is bounded
above on $\HH_z$ modulo $\Delta$.  Replacing $\beta$ by
$\beta+c_0$ adds $\Im c_0$ to $D$; a sufficiently negative imaginary
constant gives \eqref{eq:D-negative}.  Formula \eqref{eq:real-det} then
shows that $\Pi(z)$ is a real isomorphism for every $z$, so its integral
columns form a lattice in $\C^2$.
\end{proof}

Define
\[
 \mathcal T=(\HH_z\times\C^2)/\Lambda,
 \qquad
 t_\lambda(z,\zeta)=(z,\zeta+\Pi(z)\lambda).
\]
The lattice action is free and properly discontinuous.  Indeed, over a
compact subset of $\HH_z$, the real isomorphisms $\Pi(z)$ and their
inverses are uniformly bounded.  Hence, for compact subsets of
$\HH_z\times\C^2$, only finitely many lattice translations can produce
an intersection.  Thus $\mathcal T$ is a complex manifold and the
projection $\mathcal T\to\HH_z$ is a holomorphic family of compact
complex two-tori.  We use the same notation $t_\lambda$ for the
corresponding fibre loop in fundamental-group calculations.  The maps
\[
 \wt g(z,\zeta)=(gz,R_g(z)\zeta)
\]
normalize the lattice action.  Let
$\HH_z^\circ$ be the complement of the $\Delta$-orbits of $z_1,z_2$.

\begin{corollary}\label{cor:smooth-family}
The quotient
\[
 J=(\mathcal T|_{\HH_z^\circ})/\Delta
 \longrightarrow B^\circ
\]
is a proper holomorphic submersion with fibres complex two-tori and a
holomorphic zero section.  In the marking $H_1(F;\Z)=\Lambda$, the
clockwise meridians about $p_1,p_2,p_0$ have monodromies
$A_1,A_2,M_0$.
\end{corollary}

\begin{proof}
The action on $\HH_z^\circ$ is free and properly discontinuous, hence so
is the induced action on the torus family.  Over a simply connected disc
$U\Subset B^\circ$, choose one lift $\widetilde U\subset\HH_z^\circ$.
The quotient above $U$ is then the family
$(\widetilde U\times\C^2)/\Lambda$, with lattice generators depending
holomorphically on the base.  A fixed real fundamental parallelepiped for
$\Lambda\otimes\R/\Lambda$, transported by $\Pi(z)$, gives a compact
set surjecting onto the inverse image of every compact subset of $U$.
Thus the map is proper locally on the base, hence proper globally.  The
same local description shows that it is a holomorphic submersion and that
the zero section descends.

Parallel transport in the trivial marked family over $\HH_z^\circ$,
followed by the identification through $\wt g$, acts on the lattice by
$M_g^{-1}$, which is $A_j$ for $g_j$ and $M_0$ for $g_0$.
\end{proof}

\subsection{The period normal form at the cusp}
On $U_r$ put $b=\beta+\tau$.  The functions $\mu,b,h$ are holomorphic in
$t_c$, and with
\[
 C(t_c)=
 \begin{pmatrix}
 6\mu(t_c)&h(t_c)\\
 b(t_c)-h(t_c)&\mu(t_c)
 \end{pmatrix}
\]
one has
\begin{equation}\label{eq:cusp-normal-form}
 Z(z)=s(z)B_0+C(t_c\circ\pi(z)).
\end{equation}
Consequently, for $\ol\lambda\in\ol\Lambda\cong\Z^2$,
\begin{equation}\label{eq:cusp-exponential}
 \exp(2\pi iZ\ol\lambda)
 =t_c^{B_0\ol\lambda}c_{\ol\lambda}(t_c),
 \qquad
 c_{\ol\lambda}(t_c)=\exp(2\pi iC(t_c)\ol\lambda).
\end{equation}


\section{The toric filling at the cusp}

\subsection{The infinite toric threefold}
Let
\[
 N'=\Z^2\oplus\Z,
 \qquad M'=\Hom(N',\Z),
\]
and write points of $N'_\R$ as $(y,y_3)$.  Let $\mathcal T_{A_2}$ be
the triangulation of $\R^2$ with vertices $\Z^2$ and triangles
\[
 \{v,v+e_1,v+e_2\},
 \qquad
 \{v+e_1,v+e_2,v+e_1+e_2\}.
\]
For every cell $P$ let $\wh P$ be the cone over $P\times\{1\}$, and let
$\Sigma$ be the fan consisting of these cones and their faces.

\begin{lemma}[The toric model]\label{lem:toric-model}
The collection $\Sigma$ is a smooth countable fan with support
$(\R^2\times\R_{>0})\cup\{0\}$.  The associated toric space
$Y=Y_\Sigma$ is a connected Hausdorff second-countable complex
threefold.  On its dense torus $(\C^*)^3$, use coordinates
$(x_1,x_2,t_c)$, where $t_c$ is the character corresponding to
$(0,0,1)\in M'$.  If a maximal cone has vertices $v_0,v_1,v_2$ and
$m_0,m_1,m_2$ is the dual basis, then
\[
 U_\sigma=\Spec\C[\wh\sigma^\vee\cap M']\cong\C^3,
 \qquad z_i=\chi^{m_i},
\]
and
\begin{equation}\label{eq:toric-local-product}
 t_c=z_0z_1z_2.
\end{equation}
The divisor $Y_0=t_c^{-1}(0)$ is reduced and equals
$\bigcup_{v\in\Z^2}E_v$.  Each $E_v$ is the toric del Pezzo surface
$dP_6$, and its toric boundary is the hexagon of six $(-1)$-curves indexed
by the six edges through $v$.
\end{lemma}

\begin{proof}
For the lower triangle the determinant of the three primitive generators
$(v,1),(v+e_1,1),(v+e_2,1)$ is $1$; for the upper triangle it is $-1$.
Thus every cone is unimodular.  Two cones meet in the cone over the common
face of their cells, and the support is as stated because a point
$(y,y_3)$ with $y_3>0$ lies over the cell containing $y/y_3$.
Glue the affine toric charts $U_\sigma$ along the open subspaces
$U_{\sigma\cap\tau}$.  The fan condition implies that if two points
represented in $U_\sigma$ and $U_\tau$ are identified through a chain of
charts, then they are already identified in $U_{\sigma\cap\tau}$.
Consequently the union of any two charts is the toric space of the finite
fan consisting of the faces of $\sigma$ and $\tau$; it is Hausdorff by
the standard separation argument for fans \cite[Secs.~1.4,3.1]{Fulton}.
Thus the full gluing is Hausdorff.  Every chart is smooth, the atlas is
countable, and all charts meet the dense torus.  Hence $Y$ is a connected
second-countable complex threefold.

The height character $(0,0,1)$ is nonnegative on the fan and takes value
one on every ray $(v,1)$.  Hence its divisor is
$\sum_vE_v$ with multiplicity one.  In a maximal chart,
$(0,0,1)=m_0+m_1+m_2$, proving
\eqref{eq:toric-local-product} and reducedness.  The star of the ray
$(v,1)$, modulo that ray, is the complete smooth fan with rays in the six
directions
$\pm e_1,\pm e_2,\pm(e_1-e_2)$, which is the fan of $dP_6$.
\end{proof}

\subsection{The lattice action}
Identify the cocharacter lattice of $(x_1,x_2,1)$ with
$\Lambda_{\mathrm{tor}}$ by $e_1\leftrightarrow\wh w$ and
$e_2\leftrightarrow\wh\delta$.  For
$\ol\lambda\in\ol\Lambda$ define
\[
 \varphi_{\ol\lambda}(y,y_3)
 =(y+y_3B_0\ol\lambda,y_3).
\]
It preserves the fan and induces a toric automorphism $\Phi_{\ol\lambda}$
of $Y$.  On the dense torus,
\[
 \Phi_{\ol\lambda}(x,t_c)
 =(t_c^{B_0\ol\lambda}x,t_c).
\]
Using \eqref{eq:cusp-exponential}, set
\begin{equation}\label{eq:Psi-action}
 \Psi_{\ol\lambda}(x,t_c)
 =\bigl(c_{\ol\lambda}(t_c)t_c^{B_0\ol\lambda}x,t_c\bigr).
\end{equation}
The formula extends holomorphically to $Y$ because it is the product of
$\Phi_{\ol\lambda}$ with a holomorphic torus translation.  The shears
$\varphi_{\ol\lambda}$ fix the horizontal cocharacter lattice
$\Z^2\oplus0$, so the toric automorphisms commute with these translations.
Since $C(t_c)(\ol\lambda+\ol\mu)=C(t_c)\ol\lambda+C(t_c)\ol\mu$, one has
$\Psi_{\ol\lambda}\Psi_{\ol\mu}=\Psi_{\ol\lambda+\ol\mu}$.  Thus
\eqref{eq:Psi-action} defines an action of $\ol\Lambda\cong\Z^2$.
The compact horizontal torus
$T_f=\Lambda_{\mathrm{tor}}\otimes S^1$ commutes with this action and
therefore acts on the quotient constructed below, preserving $t_c$.

For a torus point with $0<|t_c|<1$ define
\[
 y(p)=\frac{\log|x(p)|}{\log|t_c(p)|}\in\R^2.
\]
Let $\theta_i^\sigma$ be the barycentric coordinates of a maximal
triangle $\sigma$, and put
\[
 \sigma(\eta)=\{y:\theta_i^\sigma(y)>-\eta,
 \ i=0,1,2\}.
\]
For $S\geq1$ write
$U_\sigma(S)=\{p\in U_\sigma:|z_i(p)|<S\}$.

\begin{lemma}[Toric chart estimates]\label{lem:chart-estimates}
For a torus point $p$,
\[
 p\in U_\sigma(S)
 \quad\Longleftrightarrow\quad
 y(p)\in\sigma\left(\frac{\log S}{|\log|t_c(p)||}\right).
\]
Every point of $|t_c|<1$ lies in the closed unit polydisc of some maximal
toric chart.  In particular the sets $U_\sigma(2)$ cover $Y_\eps$, and
every compact subset has a finite subcover.
\end{lemma}

\begin{proof}
In the coordinates of \cref{lem:toric-model},
\[
 \log|z_i(p)|=\log|t_c(p)|\,\theta_i^\sigma(y(p)),
\]
which gives the first assertion because $\log|t_c|<0$.  For a torus point,
choose a triangle containing $y(p)$.

Let now $p$ lie in the orbit associated with the cone over a cell $P$.
For $\bar m\in M'\cap\widehat P^{\perp}$, the assignment
$\bar m\mapsto\log|\chi^{\bar m}(p)|$ is a real linear functional; under
the natural pairing it is represented by a vector
$\operatorname{Log}_P(p)$ in
$N'_\R/\operatorname{span}_\R\widehat P$.  For a maximal triangle
$\sigma\supset P$, let $\bar\sigma$ be the image of its cone in this
quotient.  The cones $\bar\sigma$, as $\sigma$ ranges over the maximal
triangles containing $P$, cover the quotient vector space: if $v$ is a
quotient vector and $q_0$ lies in the relative interior of
$P\times\{1\}$, then for sufficiently small $u>0$ the height-one
normalization of $q_0+u\widetilde v$ lies in a triangle containing $P$.
Choose $\sigma$ with
$-\operatorname{Log}_P(p)\in\bar\sigma$.  If $\bar m$ is a dual
coordinate on the orbit, then
$\log|\chi^{\bar m}(p)|=\langle\bar m,\operatorname{Log}_P(p)\rangle$;
the chosen inclusion therefore makes every nonzero chart coordinate have
modulus at most one.  This proves the second assertion.  The last assertion
follows from compactness.
\end{proof}

\begin{theorem}[Cusp filling]\label{thm:cusp-filling}
For sufficiently small $\eps>0$, the action of $\ol\Lambda$ on
\[
 Y_\eps=\{|t_c|<\eps\}\subset Y
\]
is free and properly discontinuous.  The quotient
\[
 N_0=Y_\eps/\ol\Lambda
\]
is a connected Hausdorff complex threefold, and $t_c$ descends to a proper
surjective holomorphic map
\[
 f_0:N_0\longrightarrow\Delta_0=\{|t_c|<\eps\}.
\]
It is a submersion away from $0$, and its nonzero fibres are the tori of
\cref{cor:smooth-family}.  The central fibre
\[
 W=f_0^{-1}(0)
\]
is reduced, irreducible and has normal crossings.  Its normalization is
$dP_6$; the normalization map identifies opposite sides of the
anticanonical hexagon.  The double locus consists of three rational curves
through two triple points.

Moreover
\[
 \pi_1(N_0)\cong\ol\Lambda,
\]
the inclusion of a nearby fibre sends $\Lambda$ to
$\ol\Lambda=\Lambda/\Lambda_{\mathrm{tor}}$, and the loop obtained by
holding $(x_1,x_2)=(1,1)$ and circling $t_c=0$ is null-homotopic in $N_0$.
\end{theorem}

\begin{proof}
Put
\[
 R(t_c)=-2\pi\Im C(t_c),
 \qquad
 B_{t_c}=B_0+\frac{R(t_c)}{\log|t_c|}.
\]
Choose $\eps$ so small that
$\|B_{t_c}-B_0\|\leq\frac12$ for $0<|t_c|<\eps$.
Since $B_0$ is orthogonal in the chosen bases,
\begin{equation}\label{eq:Bt-bound}
 |B_{t_c}\ol\lambda|\geq\frac12|\ol\lambda|.
\end{equation}

\smallskip
\noindent\emph{Freeness.}
On the dense torus, a fixed point of $\Psi_{\ol\lambda}$ with
$\ol\lambda\ne0$ would satisfy
\[
 \log|t_c|\,B_0\ol\lambda+R(t_c)\ol\lambda=0,
\]
which contradicts the choice of $\eps$.  On $Y_0$, the action sends the
orbit attached to a cell $P$ to the orbit attached to
$P+B_0\ol\lambda$.  No nonzero translation preserves a bounded cell, so
there are no fixed points there either.

\smallskip
\noindent\emph{Proper discontinuity.}
Equation \eqref{eq:Psi-action} gives, on the dense torus,
\begin{equation}\label{eq:position-translation}
 y(\Psi_{\ol\lambda}p)=y(p)+B_{t_c(p)}\ol\lambda.
\end{equation}
Fix maximal triangles $\sigma,\sigma'$.  By
\cref{lem:chart-estimates}, the positions of torus points in
$U_\sigma(2)\cap Y_\eps$ and $U_{\sigma'}(2)\cap Y_\eps$ lie in
bounded sets $\sigma(\eta_0)$ and $\sigma'(\eta_0)$, where
$\eta_0=\log2/|\log\eps|$.  If
$\Psi_{\ol\lambda}U_\sigma(2)$ meets $U_{\sigma'}(2)$, density of
the torus in every chart supplies a torus point in the intersection.
Then
\[
 \frac12|\ol\lambda|
 \leq |B_{t_c}\ol\lambda|
 \leq
 \sup\{|a-b|:a\in\sigma'(\eta_0),\ b\in\sigma(\eta_0)\},
\]
so only finitely many $\ol\lambda$ are possible.  A compact subset is
covered by finitely many $U_\sigma(2)$; therefore only finitely many of
its translates meet it.  This is proper discontinuity.

The orbit space is Hausdorff and the orbit map is a covering.  If $K$ is
a compact neighbourhood of a point $p$, proper discontinuity leaves only
finitely many nontrivial translates of $K$ that meet $K$; freeness and
Hausdorffness allow $K$ to be shrunk to an open neighbourhood disjoint
from all its nontrivial translates.  The same finite-translate argument
separates two distinct orbits.  Since all deck transformations are
biholomorphic, the quotient is a complex manifold.  The set $Y_\eps$ is connected: in the coordinates of a maximal chart,
$U_\sigma\cap Y_\eps=\{|z_0z_1z_2|<\eps\}$ is star-shaped and meets
the connected dense set
$(\C^*)^2\times\{0<|t_c|<\eps\}$.  It is second countable, and the open
quotient map shows that $N_0$ is second countable as well.  Hence $N_0$ is
a connected complex threefold.

\smallskip
\noindent\emph{Properness.}
Fix $0<\eps_1<\eps$.  The matrices $B_{t_c}$ and their inverses are
uniformly bounded for $0<|t_c|\leq\eps_1$.  Given a torus point $p$,
put $u=B_{t_c}^{-1}y(p)$ and choose $\ol\lambda\in\Z^2$ with
$|u+\ol\lambda|_\infty\leq1/2$.  Then
\[
 |y(\Psi_{\ol\lambda}p)|
 =|B_{t_c}(u+\ol\lambda)|
 \leq C
\]
for a constant independent of $p$.  Choose the finite set of maximal
triangles meeting the closed ball $\overline{B}(0,C+2)$ and, in each
corresponding chart, the compact unit polydisc
$\{|z_i|\leq1,\ |t_c|\leq\eps_1\}$.  By
\cref{lem:chart-estimates}, every translated torus point above the closed
disc lies in their union.

For a point of $Y_0$ in the orbit attached to a cell $P$, choose a vertex
$v$ of $P$.  Unimodularity of $B_0$ gives $\ol\lambda$ with
$B_0\ol\lambda=-v$, so the translated cell contains $0$ and is a face
of one of the finitely many triangles adjacent to $0$.  Enlarging the
preceding compact union by the corresponding closed chart polydiscs gives
a compact set $K_0$ meeting every orbit over $|t_c|\leq\eps_1$.
Therefore
$f_0^{-1}(|t_c|\leq\eps_1)$ is the compact image of $K_0$.  Any compact
subset of the open disc lies in such a closed subdisc, proving that $f_0$
is proper.

\smallskip
\noindent\emph{The fibres and local structure.}
For $t_c\ne0$, choose $s$ with $e^{2\pi is}=t_c$.  The exponential map
$\C^2\to(\C^*)^2$ identifies the quotient fibre with
\[
 \C^2/\bigl(\Z^2+(sB_0+C(t_c))\Z^2\bigr)
 =\C^2/\Pi(z)\Lambda,
\]
where $s=s(z)$; this is the torus family of
\cref{cor:smooth-family}.  These fibres exist
for every $0<|t_c|<\eps$, while $Y_0$ is nonempty, so $f_0$ is
surjective.  On $t_c\ne0$ the character $t_c:(\C^*)^3\to\C^*$ is a
submersion, and the quotient map is locally biholomorphic.  Hence $f_0$
is a submersion away from the origin.  Finally, suppose that exactly the coordinates indexed by a nonempty set
$I\subset\{0,1,2\}$ vanish at a chosen point.  The product of the remaining
coordinates is a holomorphic unit; absorbing that unit into one coordinate
indexed by $I$, \eqref{eq:toric-local-product} gives local coordinates on
$N_0$ in which $f_0$ is $z_0$, $z_0z_1$, or $z_0z_1z_2$, according as
$|I|=1,2,3$.  Thus $W$ is reduced and has normal crossings.

Because $B_0\ol\Lambda=\Z^2$, the action on the height-one lattice
is the full translation action.  Every cell has a vertex, so every toric
orbit in $Y_0$ is equivalent to an orbit contained in $E_0$; hence the map
$E_0\to W$ is surjective.  It is injective on the open torus orbit.  Two
orbits in $E_0$ can be identified only when cells $P,P'$ containing $0$
satisfy $P'=P+v$ for a nonzero integral translation.  For edges this
forces $P=\{0,v\}$ and $P'=\{0,-v\}$.  No nonzero translation preserves
an edge, so each boundary curve maps isomorphically to its image.  Thus
the six boundary curves are identified in the three opposite pairs and no
other curve identifications occur.  For triangles, the six triangles through $0$ fall into two
translation orbits, alternating around the hexagon; their images are the
two triple points.  The fibres of $E_0\to W$ are therefore finite.  The map is proper because
$E_0$ is compact and $W$ is Hausdorff; a proper holomorphic map with finite
fibres is finite.  Since $E_0$ is irreducible and maps onto $W$, the surface
$W$ is irreducible.  The finite map $E_0\to W$ is generically one-to-one,
with normal source and reduced target; it is therefore the normalization.  This proves the stated
description of the double and triple loci.

\smallskip
\noindent\emph{Fundamental group.}
The toric threefold $Y$ is simply connected.  For $n\ge1$, let
$\Sigma_n$ be the finite subfan consisting of the cones over the cells of
$\mathcal T_{A_2}$ contained in the square $[-n,n]^2$, together with all
their faces.  The corresponding toric spaces $Y_{\Sigma_n}$ form an
increasing open cover of $Y$.  Each contains the rays through
$(0,0,1),(1,0,1),(0,1,1)$, which generate $N'$ over $\Z$.  The standard
toric fundamental-group formula \cite[Sec.~3.2]{Fulton} therefore gives
$\pi_1(Y_{\Sigma_n})=0$.  Every loop has compact image and hence lies
in some $Y_{\Sigma_n}$, so $\pi_1(Y)=0$.

For $0<u\leq1$, the torus element $(1,1,u)\in(\C^*)^3$ acts on $Y$ and
satisfies $t_c((1,1,u)\cdot p)=u\,t_c(p)$.  A loop in $Y_\eps$ bounds a
disc in $Y$; for $u>0$ sufficiently small, translating that disc by
$(1,1,u)$ places it in $Y_\eps$.  On the boundary, the path of torus
elements $(1,1,u_s)$, with $u_s$ decreasing from $1$ to $u$, gives a
homotopy inside $Y_\eps$ from the original loop to the translated one.
Hence $Y_\eps$ is simply connected.
Therefore the covering $Y_\eps\to N_0$ is universal and
$\pi_1(N_0)=\ol\Lambda$.

On a nonzero fibre the covering $(\C^*)^2\to f_0^{-1}(t_c)$ has deck group
$\ol\Lambda$, so the induced map on fundamental groups is the projection
$\Lambda\to\ol\Lambda$.  Finally, the curve
$\{(1,1,t_c):|t_c|<\eps\}$ is a holomorphic disc in a toric chart; its
boundary is the stated meridian, which is therefore null-homotopic.
\end{proof}


\begin{proposition}[Identification over the punctured cusp disc]
\label{prop:cusp-identification}
After decreasing $\eps$ if necessary, put
$U_\eps=U_r\cap\pi^{-1}(D_0^*)$.  The map
\begin{equation}\label{eq:cusp-exp-map}
 E_0^{\exp}:U_\eps\times\C^2\longrightarrow
 Y_\eps\cap(\C^*)^3,
 \qquad
 E_0^{\exp}(z,\zeta)=
 \bigl(e^{2\pi i\zeta_1},e^{2\pi i\zeta_2},e^{2\pi is(z)}\bigr),
\end{equation}
induces a biholomorphism
\begin{equation}\label{eq:G0}
 G_0:J|_{D_0^*}\xrightarrow{\sim}N_0\setminus W,
 \qquad D_0^*=\{0<|t_c|<\eps\}.
\end{equation}
It identifies the marking $H_1(F;\Z)=\Lambda$ on both sides and carries the
zero section to the section represented by $(x_1,x_2)=(1,1)$.  The latter
extends across $t_c=0$ as a holomorphic section of $f_0$, meeting $W$ in
its smooth locus.
\end{proposition}

\begin{proof}
Let $\wt g_0(z,\zeta)=(g_0z,\zeta)$; the latter
formula follows from $R_0=I_2$.  Let $H$ be the subgroup generated by
$\wt g_0$ and the translations $t_\lambda$ with
$\lambda\in\Lambda_{\mathrm{tor}}$.  It is normal in
$\Lambda\rtimes\langle g_0\rangle$: conjugation by $\wt g_0$ preserves
the translations by $\Lambda_{\mathrm{tor}}$, while
\[
 t_\mu\,\wt g_0\,t_\mu^{-1}
 =t_{\mu-M_0\mu}\,\wt g_0\in H
 \qquad(\mu\in\Lambda),
\]
because $(M_0-I)\Lambda\subset\Lambda_{\mathrm{tor}}$.

The map $s:U_\eps\to\{\Im s>c\}$ is a biholomorphism and satisfies
$s\circ g_0=s-1$.  Hence the last coordinate in
\eqref{eq:cusp-exp-map} is the universal covering of the punctured disc,
with deck group generated by $g_0$.  The first two coordinates are the
ordinary exponential covering $\C^2\to(\C^*)^2$, whose deck lattice is
$\Z^2=\Pi(z)\Lambda_{\mathrm{tor}}$.  It follows that
$E_0^{\exp}$ is a surjective local biholomorphism and that its fibres are
exactly the $H$-orbits.

Write $\lambda=\lambda_{\mathrm{tor}}+\lambda'$ and let
$\ol\lambda$ be its image in $\ol\Lambda$.  By
\eqref{eq:cusp-normal-form} and \eqref{eq:cusp-exponential},
\[
 E_0^{\exp}\bigl(z,\zeta+\Pi(z)\lambda\bigr)
 =\Psi_{\ol\lambda}\bigl(E_0^{\exp}(z,\zeta)\bigr).
\]
Thus $E_0^{\exp}$ identifies $(U_\eps\times\C^2)/H$ with
$Y_\eps\cap(\C^*)^3$ and intertwines the residual action
$(\Lambda\rtimes\langle g_0\rangle)/H\cong\ol\Lambda$ with the action
$\Psi$ of \eqref{eq:Psi-action}.  Since $U_\eps$ is the only component of
$\pi^{-1}(D_0^*)$ stabilized by $g_0$ and its other $\Delta$-translates
are disjoint, the quotient by the full residual action is precisely
\eqref{eq:G0}.

The construction sends the cycles in $\Lambda_{\mathrm{tor}}$ to the
fundamental cycles of $(\C^*)^2$ and the quotient
$\Lambda/\Lambda_{\mathrm{tor}}$ to the deck group of a fibre of
$N_0\setminus W$; hence it preserves the marking.  Finally
$E_0^{\exp}(z,0)=(1,1,t_c)$.  In the toric chart corresponding to the
triangle with vertices $0,e_1,e_2$, the coordinates are
$(t_c/(x_1x_2),x_1,x_2)$, so
\[
 \{(1,1,t_c):|t_c|<\eps\}
\]
is a holomorphic disc meeting the divisor $t_c=0$ transversely at a point
of the open orbit of $E_0$.  Since every $\Psi_{\ol\lambda}$ preserves $t_c$ and the disc contains
exactly one point over each value of $t_c$, the quotient map is injective
on this disc.  Its image is therefore the required holomorphic section.
\end{proof}

\begin{proposition}[Boundary of the cusp piece]\label{prop:cusp-boundary}
For $0<\rho<\eps$, the boundary
$\mathcal M_0=f_0^{-1}(|t_c|=\rho)$ is the mapping torus of the monodromy
$M_0$ on $F$.  With $a_0$ the toric meridian,
\[
 \pi_1(\mathcal M_0)=\Lambda\rtimes_{M_0}\Z\langle a_0\rangle.
\]
The inclusion $\mathcal M_0\hookrightarrow N_0$ sends
$\lambda\mapsto\ol\lambda\in\ol\Lambda$ and $a_0\mapsto1$; its
kernel is the normal closure of
$\Lambda_{\mathrm{tor}}\cup\{a_0\}$.
\end{proposition}

\begin{proof}
Ehresmann's theorem identifies the restriction over the circle with the
mapping torus of the geometric monodromy, whose translation part is
isotopic to zero and whose linear part is $M_0$.  The map on the fibre
fundamental group is the projection established in
\cref{thm:cusp-filling}, and the meridian is null-homotopic there.  The
resulting quotient of the displayed semidirect product is
$\Lambda/\Lambda_{\mathrm{tor}}=\ol\Lambda$, which is
$\pi_1(N_0)$; hence the stated normal subgroup is exactly the kernel.
\end{proof}


\section{The finite fillings and the compact complex threefold}

\subsection{Free logarithmic transforms}
For $j=1,2$, choose a $g_j$-invariant disc $\Delta_j\ni z_j$ and a
coordinate $s_j$ such that
\[
 s_j(z_j)=0,
 \qquad
 s_j\circ g_j=e^{-2\pi i/m_j}s_j.
\]
The quotient coordinate $t_j$ on $D_j=\Delta_j/\langle g_j\rangle$ is
defined by $t_j\circ\pi=s_j^{m_j}$.  Let
\[
 J'_j=(\Delta_j\times\C^2)/\Lambda
\]
be the restriction of the marked torus family.  Define
\begin{equation}\label{eq:log-action}
 \wt g_j^{\mathrm{log}}(z,\zeta)
 =\left(g_jz,R_j(z)\zeta+\frac1{m_j}\Pi(g_jz)v_j\right).
\end{equation}
The identity
$R_j(z)\Pi(z)=\Pi(g_jz)A_j$ implies
\[
 \wt g_j^{\mathrm{log}}\,t_\lambda
 =t_{A_j\lambda}\,\wt g_j^{\mathrm{log}},
\]
so \eqref{eq:log-action} descends to $J'_j$.  In the smooth
trivialization
\[
 \Delta_j\times(\Lambda\otimes\R)/\Lambda\longrightarrow J'_j,
 \qquad (z,x)\longmapsto[(z,\Pi(z)x)],
\]
it is
\begin{equation}\label{eq:log-action-flat}
 (z,x)\longmapsto(g_jz,A_jx+v_j/m_j).
\end{equation}

\begin{lemma}[Fixed points of an affine torus map]\label{lem:affine-fixed}
Let $A\in\GL(\Lambda)$ have finite order and let $b\in\Lambda\otimes\Q$.
The affine map $x\mapsto Ax+b$ of $(\Lambda\otimes\R)/\Lambda$ has a fixed
point if and only if
\[
 \phi(b)\in\Z
 \quad\text{for every }\phi\in V\text{ satisfying }\phi\circ A=\phi.
\]
\end{lemma}

\begin{proof}
A fixed point exists exactly when
$b\in(A-I)(\Lambda\otimes\R)+\Lambda$.  Necessity is immediate.  For the
converse, average the powers of $A$ to obtain the rational projection
$p$ onto $\ker(A-I)_\R$ with kernel $\im(A-I)_\R$.  The lattice
$p(\Lambda)$ is dual to the invariant lattice
$\{\phi\in V:\phi A=\phi\}$.  The integrality assumption therefore gives
$pb=p\lambda$ for some $\lambda\in\Lambda$, so
$b-\lambda\in\im(A-I)_\R$.
\end{proof}

\begin{theorem}[The two logarithmic fillings]\label{thm:finite-fillings}
The map \eqref{eq:log-action} generates a free cyclic action of order
$m_j$ on $J'_j$.  Its quotient
\[
 N_j=J'_j/\langle\wt g_j^{\mathrm{log}}\rangle
\]
is a smooth complex threefold, and $s_j^{m_j}$ descends to a proper
holomorphic map
\[
 f_j:N_j\longrightarrow D_j
\]
with divisor
\[
 f_j^*(p_j)=m_jS_j,
\]
where $S_j$ is a smooth compact surface.

Over the punctured disc there is a biholomorphism
\[
 G_j:N_j\setminus S_j\longrightarrow J|_{D_j^*}.
\]
Under this biholomorphism the zero section of $J$ is represented on the
cover by
\begin{equation}\label{eq:zero-section-log}
 \varsigma_j(z)=-\frac{\log s_j(z)}{2\pi i}\Pi(z)v_j.
\end{equation}
Following this section around $s_j$ once counterclockwise changes its
lift by $-\Pi(z)v_j$.
\end{theorem}

\begin{proof}
Because $A_jv_j=v_j$, the $m_j$-th power of
\eqref{eq:log-action-flat} is translation by $v_j\in\Lambda$, hence the
identity.  A nontrivial power can have a fixed point only over $z_j$.
A direct calculation gives
\[
 V^{T_1}=V^{T_1^2}=\langle\gamma,\psi_1\rangle,
 \qquad \psi_1=2u+w+3\delta,
\]
and
\[
 V^{T_2}=V^{T_2^2}=V^{T_2^3}=\langle\gamma,\psi_2\rangle,
 \qquad \psi_2=u+w+2\delta.
\]
The $k$-th power of \eqref{eq:log-action-flat} has translation part
$kv_j/m_j$.  For $v_1=\eps$ one has
$\gamma(v_1)=1$ and $\psi_1(v_1)=0$; hence for $k=1,2$ the value
$k\gamma(v_1)/3$ is not integral.  For $v_2=-\eps'$ one has
$\gamma(v_2)=-1$ and $\psi_2(v_2)=0$; hence for $k=1,2,3$ the value
$-k/4$ is not integral.  The fixed-point criterion in
\cref{lem:affine-fixed} therefore excludes fixed points for every
nontrivial power.

A free finite holomorphic action has a smooth quotient.  The function
$s_j^{m_j}$ is invariant, properness follows from properness of
$J'_j\to\Delta_j$ and finiteness of the quotient, and its divisor is
$m_j$ times the reduced central fibre.

On $\Delta_j^*$ define the section
\[
 \sigma_j(z)=\frac{\log s_j(z)}{2\pi i}\Pi(z)v_j.
\]
Changing the logarithm changes $\sigma_j$ by a period, so it is well
defined on the torus family.  Translation by $\sigma_j$ conjugates
$\wt g_j^{\mathrm{log}}$ to the linear deck transformation $\wt g_j$.
Indeed
$s_j(g_jz)=e^{-2\pi i/m_j}s_j(z)$ and
$R_j(z)\Pi(z)v_j=\Pi(g_jz)A_jv_j=\Pi(g_jz)v_j$, so the translation terms
cancel.  This gives $G_j$.  The zero section in $N_j$-coordinates is
$-\sigma_j$, proving \eqref{eq:zero-section-log} and the stated drift.
\end{proof}

\begin{proposition}[Topology of the logarithmic pieces]\label{prop:pi1-Nj}
The radial contraction of the $s_j$-disc is equivariant and induces a
deformation retraction $N_j\simeq S_j$.  Its fundamental group has the
presentation
\[
 \pi_1(N_j)=
 \left\langle \Lambda,\wh g_j\ \middle|\
 \wh g_j\lambda\wh g_j^{-1}=A_j\lambda,
 \ \wh g_j^{m_j}=t_{v_j}
 \right\rangle.
\]
For the boundary mapping torus $M_j=f_j^{-1}(|t_j|=\rho)$ one has
$\pi_1(M_j)=\Lambda\rtimes_{A_j}\Z\langle\wh g_j\rangle$, and the
kernel of $\pi_1(M_j)\to\pi_1(N_j)$ is normally generated by
$\wh g_j^{m_j}t_{-v_j}$.
\end{proposition}

\begin{proof}
Before quotienting, the local model is smoothly
$A_0^{(j)}\times\Delta_{s_j}$, where
$A_0^{(j)}=\C^2/L(z_j)$, and the generator acts by
\[
 (x,s_j)\longmapsto(A_jx+v_j/m_j,
 e^{-2\pi i/m_j}s_j).
\]
Radial contraction of the disc is equivariant.  The universal cover of
$S_j$ is $\Lambda\otimes\R$, and the group generated by translations
$t_\lambda$ and the affine lift
$x\mapsto A_jx+v_j/m_j$ has the displayed relations; its $m_j$-th power
is $t_{v_j}$.  On the boundary, the lift of one $m_j$-th of the clockwise
base circle is an infinite-order element $\wh g_j$.  Traversing it
$m_j$ times gives the lattice translation $v_j$ together with the full
boundary circle of the covering disc.  Killing that boundary circle on
filling the disc gives exactly the normal relation
$\wh g_j^{m_j}t_{-v_j}$.
\end{proof}

\subsection{Gluing}
Shrink $D_0,D_1,D_2$ so that they are pairwise disjoint and
\[
 B=B^\circ\cup D_0\cup D_1\cup D_2.
\]
Use the biholomorphism $G_0$ of \cref{prop:cusp-identification} and the
biholomorphisms $G_1,G_2$ of \cref{thm:finite-fillings}.

Define
\[
 X=(J\sqcup N_0\sqcup N_1\sqcup N_2)/\sim,
\]
where points in $J|_{D_0^*}$ are identified through $G_0$, and points in
$N_j\setminus S_j$ through $G_j$.

\begin{theorem}[The compact complex threefold]\label{thm:global-X}
The space $X$ is a connected compact Hausdorff complex threefold.  The
local maps glue to a proper surjective holomorphic map
\[
 f:X\longrightarrow B=\PP^1
\]
with connected fibres.  Over $B^\circ$ its fibres are complex two-tori;
$f^{-1}(p_0)=W$ is the reduced normal crossings surface of
\cref{thm:cusp-filling}; and
\[
 f^*(p_1)=3S_1,
 \qquad
 f^*(p_2)=4S_2.
\]
\end{theorem}

\begin{proof}
The three overlap regions in $J$ are disjoint.  Hence every nonsingleton
equivalence class consists of exactly two points, and the quotient maps
from the four pieces are open embeddings.  Their transition maps are the
biholomorphisms $G_0^{\pm1},G_1^{\pm1},G_2^{\pm1}$, so they define a
complex atlas on $X$.  Since the four pieces are second countable and their
images are open, $X$ is second countable.

Points with different images in $B$ are separated by inverse images of
disjoint base neighbourhoods.  Points over the same base point lie in one
of the Hausdorff local models, proving that $X$ is Hausdorff.  Properness
is local on the base and holds for each of the four local maps; therefore
$f$ is proper and $X=f^{-1}(B)$ is compact.  Connectedness and the fibre
descriptions follow from the corresponding properties of the local
models.
\end{proof}


\section{The fundamental group}
Orient every meridian clockwise.  Let $\rho_j$ be the lift of the meridian
about $p_j$ along the zero section of $J$.  Since $B^\circ$ is a pair of
pants,
\[
 \pi_1(J)=\Lambda\rtimes F(\rho_1,\rho_2),
 \qquad
 \rho_j\lambda\rho_j^{-1}=A_j\lambda,
 \qquad
 \rho_0=(\rho_1\rho_2)^{-1}.
\]

\begin{lemma}[Relations imposed by the fillings]
\label{lem:filling-relations}
Under the gluing maps, the clockwise meridian lifts satisfy
\begin{equation}\label{eq:finite-pi1-relations}
 \rho_1^3=t_{v_1},
 \qquad
 \rho_2^4=t_{v_2},
\end{equation}
and the cusp filling kills $\Lambda_{\mathrm{tor}}$ and imposes
\begin{equation}\label{eq:cusp-pi1-relations}
 \rho_1\rho_2=1.
\end{equation}
\end{lemma}

\begin{proof}
By \cref{prop:pi1-Nj}, attaching $N_j$ kills
$\wh g_j^{m_j}t_{-v_j}$.  A clockwise meridian in the $t_j$-disc lifts
through $t_j=s_j^{m_j}$ to a path on which $s_j$ turns through
$-2\pi/m_j$.  Its endpoint is obtained by the deck transformation
$g_j$, because $s_j\circ g_j=e^{-2\pi i/m_j}s_j$.  In the coordinates
of $N_j$, the lift of the zero section is
$\varsigma_j=-\log(s_j)\Pi v_j/(2\pi i)$.  The identity
\[
 R_j(z)\varsigma_j(z)+\frac1{m_j}\Pi(g_jz)v_j
 =\varsigma_j(g_jz)
\]
follows from $\log(s_j\circ g_j)=\log s_j-2\pi i/m_j$ and
$A_jv_j=v_j$.  Thus the gluing identifies $\rho_j$ with $\wh g_j$, with
the same sign for $j=1,2$, and
\eqref{eq:finite-pi1-relations} follows.

By \cref{prop:cusp-boundary}, attaching $N_0$ kills the normal closure of
$\Lambda_{\mathrm{tor}}$ and of the toric meridian.  Under
\cref{prop:cusp-identification}, the zero section is the disc
$(x_1,x_2)=(1,1)$, so its clockwise boundary is the toric meridian and is
also the lift $\rho_0$.  Since $\rho_0=(\rho_1\rho_2)^{-1}$,
\eqref{eq:cusp-pi1-relations} follows.
\end{proof}

\begin{theorem}\label{thm:pi1}
The compact complex threefold $X$ is simply connected.
\end{theorem}

\begin{proof}
Apply van Kampen's theorem to the open cover by the images of $J$ and
slightly thickened cusp and logarithmic pieces.  Their intersections are
connected torus bundles over annuli, so \cref{lem:filling-relations}
gives the quotient of $\pi_1(J)$ by exactly the stated normal relations.
The normal closure of $\Lambda_{\mathrm{tor}}$ under $A_1,A_2$ is
$\ker\gamma$ by \cref{prop:linear-algebra}.  The surviving image of
$\Lambda$ is therefore cyclic, generated by the image $c$ of
$\wh\gamma$, and is central because $A_1,A_2$ act trivially on the
coinvariant quotient.  By \cref{lem:filling-relations}, with
$x=\rho_1$ and $y=\rho_2$, the group has the presentation
\[
 \langle c,x,y\mid c\text{ central},\ xy=1,\ x^3=c,\ y^4=c^{-1}\rangle,
\]
because $\gamma(v_1)=1$ and $\gamma(v_2)=-1$.  Hence $y=x^{-1}$ and
$c=x^3$, while $y^4=c^{-1}$ gives $x^{-4}=x^{-3}$.  Thus $x=1$, and then
$y=c=1$.
\end{proof}


\section{Integral cohomology}
Let $F$ be a smooth fibre, marked so that
\[
 H_1(F;\Z)=\Lambda,
 \qquad
 H^1(F;\Z)=V,
 \qquad
 H^q(F;\Z)=\bigwedge^qV.
\]
Let $j:B^\circ\hookrightarrow B$ and let $L^q$ be the local system
$R^qf_*\Z|_{B^\circ}$.

\subsection{Specialization at the multiple fibres}
Put
\[
 \psi_1=2u+w+3\delta,
 \qquad
 \psi_2=u+w+2\delta.
\]
The central torus in the covering model of $N_j$ maps to $S_j$ by an
$m_j$-sheeted covering.  Using the smooth marking over $\Delta_j$, identify
that torus with a nearby fibre $F$ and denote the covering by
$\pi_j:F\to S_j$.

\begin{proposition}\label{prop:finite-specialization}
The groups $H^q(S_j;\Z)$ are torsion-free of ranks
\[
 (1,2,2,2,1),
\]
and the pullback from $S_j$ to a nearby torus is injective.  Its image in
the monodromy invariants has the following indices:
\[
\begin{array}{c|ccccc}
q&0&1&2&3&4\\ \hline
j=1&1&3&1&1&3\\
j=2&1&4&2&2&4.
\end{array}
\]
More precisely,
\[
 \pi_1^*H^1(S_1)=\langle3\gamma,\psi_1\rangle,
 \qquad
 \pi_2^*H^1(S_2)=\langle4\gamma,\psi_2\rangle.
\]
The degree-two cokernel at $p_2$ is generated by the class of
$q=uw+6\gamma\delta$, and the degree-three cokernel at $p_2$ by the class
of $\gamma uw$.
\end{proposition}

\begin{proof}
The fundamental group of $S_j$ has the presentation
\[
 \langle\Lambda,\wh g_j\mid
 \wh g_j\lambda\wh g_j^{-1}=A_j\lambda,
 \ \wh g_j^{m_j}=t_{v_j}\rangle.
\]
The matrices in \cref{prop:linear-algebra} give
\[
 (A_1-I)\Lambda
 =\{\lambda:\gamma(\lambda)=0,
   \ \psi_1(\lambda)=0\},
 \qquad
 (A_2-I)\Lambda
 =\{\lambda:\gamma(\lambda)=0,
   \ \psi_2(\lambda)=0\}.
\]
Indeed the generators of $(A_1-I)\Lambda$ are
$6\wh u-6\wh w-2\wh\delta$,
$-\wh u-\wh w+\wh\delta$ and
$\wh u-2\wh w$, which generate the displayed kernel; the analogous
statement follows from the three generators of $(A_2-I)\Lambda$.
The maps $(\gamma,\psi_j):\Lambda\to\Z^2$ are surjective, so
$\Lambda/(A_j-I)\Lambda\cong\Z^2$.

After abelianization the relation is
\[
 m_j\wh g_j=(\gamma(v_j),\psi_j(v_j)).
\]
For $v_1=\eps$ this is $(1,0)$ and for $v_2=-\eps'$ it is $(-1,0)$.
The relation vector in $\Z^3$ is therefore primitive, and
$H_1(S_j;\Z)\cong\Z^2$ in both cases.  Since $S_j$ is a free finite
quotient of a four-torus, its Euler characteristic is zero.  Thus
$b_1=b_3=2$ and $b_2=2$.  Poincaré duality gives
$H^3(S_j;\Z)\cong H_1(S_j;\Z)\cong\Z^2$.  In the universal coefficient
sequence for $H^3$, the torsion group
$\operatorname{Ext}(H_2(S_j;\Z),\Z)$ injects into this free group; hence
$H_2(S_j;\Z)$ is torsion-free.  Poincaré duality and the universal
coefficient theorem now show that every $H^q(S_j;\Z)$ is torsion-free,
with ranks $(1,2,2,2,1)$.  The cohomological transfer of $\pi_j$ satisfies
$\operatorname{tr}\circ\pi_j^*=m_j$; hence $\pi_j^*$ is injective in every
degree.

A homomorphism $\pi_1(S_j)\to\Z$ restricting to
$\xi\in V^{T_j}$ exists exactly when $m_j$ divides $\xi(v_j)$.  Since
\[
 V^{T_1}=\langle\gamma,\psi_1\rangle,
 \qquad
 V^{T_2}=\langle\gamma,\psi_2\rangle,
\]
and $\psi_j(v_j)=0$, the degree-one images are the displayed lattices.

The invariant degree-two lattices are
\[
 (\textstyle\bigwedge^2V)^{T_1}=\langle\gamma\psi_1,q\rangle,
 \qquad
 (\textstyle\bigwedge^2V)^{T_2}=\langle\gamma\psi_2,q\rangle.
\]
Their intersection matrices, with values in $\Z\vol$, are
\[
 \begin{pmatrix}0&3\\3&12\end{pmatrix},
 \qquad
 \begin{pmatrix}0&2\\2&12\end{pmatrix},
\]
with determinants $-9$ and $-4$.  The intersection form on
$H^2(S_j;\Z)$ is unimodular.  It is even.  Indeed, the covering torus has
a nowhere-vanishing holomorphic two-form $\omega_0$.  At the central
point the linear part of the deck generator is $R_j(z_j)$, so
$\omega_0$ is multiplied by $\det R_j(z_j)$.  From
\cref{prop:period-equivariance,eq:mu-values},
\[
 \det R_1(z_1)=-\rho^{-1}=e^{2\pi i/3},
 \qquad
 \det R_2(z_2)=i^{-1}=-i.
\]
The form $\omega_0$ identifies the pullback of $K_{S_j}$ with the
trivial line bundle equipped with the character
$\wh g_j\mapsto\det R_j(z_j)^{-1}$.  This character has finite order, so
$K_{S_j}$ is torsion and hence numerically trivial.  Wu's formula gives
$x^2\equiv x\cdot c_1(K_{S_j})\equiv0\pmod2$ for every
$x\in H^2(S_j;\Z)$.  Pullback multiplies the intersection form by the
covering degree $m_j$.
If $k_j$ is the index of its image in the invariant lattice, determinant
comparison gives
\[
 k_1^2=\frac{3^2}{9}=1,
 \qquad
 k_2^2=\frac{4^2}{4}=4.
\]
Thus the degree-two indices are $1$ and $2$.  The vector $q$ is not in the
image for $j=2$: otherwise $q=\pi_2^*a$ would give
$12=q^2=4a^2$, so $a^2=3$, contradicting evenness.  Hence its class
generates the cokernel.

The degree-three indices follow from the perfect pairing
$H^1\times H^3\to\Z$.  Direct calculation gives
\[
 (\textstyle\bigwedge^3V)^{T_1}
 =\langle\gamma uw,-uw\delta+2\gamma w\delta+4\gamma u\delta\rangle,
\]
\[
 (\textstyle\bigwedge^3V)^{T_2}
 =\langle\gamma uw,-uw\delta+3\gamma w\delta+3\gamma u\delta\rangle.
\]
The determinants of the pairings with the degree-one invariant lattices
are $3$ and $2$.  Let $k_j^{(r)}$ be the index of the pullback image in
degree $r$.  If bases of the pullback lattices are expressed in bases of
the invariant lattices, the determinants of the two change-of-basis
matrices are $k_j^{(1)}$ and $k_j^{(3)}$.  On the other hand, pullback
multiplies the perfect pairing on $S_j$ by $m_j$, so the determinant of
the restricted rank-two pairing is $m_j^2$.  Therefore
\[
 k_j^{(1)}k_j^{(3)}
 \,|\det( (\textstyle\bigwedge^1V)^{T_j}
            \times(\textstyle\bigwedge^3V)^{T_j})|
 =m_j^2.
\]
Using $k_1^{(1)}=3$, $k_2^{(1)}=4$ gives
$k_1^{(3)}=1$, $k_2^{(3)}=2$.  Moreover $\psi_2$ lies in the degree-one
pullback image, whereas every pairing between pullback classes in degrees
one and three is divisible by $4$.  Since
$\langle\psi_2,\gamma uw\rangle=-2$, the class $\gamma uw$ is not in
the degree-three pullback image and therefore generates its order-two
cokernel.  In degree four pullback is multiplication by the covering
degree.
\end{proof}

\subsection{Specialization at the cusp}
\begin{theorem}[Integral invariant cycles at the cusp]\label{thm:cusp-specialization}
For every $q$, the specialization map
\[
 \operatorname{sp}^q:H^q(W;\Z)\longrightarrow H^q(F;\Z)
\]
is injective and has image
\[
 H^q(F;\Z)^{T_0}.
\]
In particular
\[
 H^q(W;\Z)\cong
 \Z,\ \Z^2,\ \Z^4,\ \Z^2,\ \Z
 \qquad(q=0,1,2,3,4),
\]
and all these groups are torsion-free.
\end{theorem}

\subsection{The integral direct-image sheaves}
Write
\[
 R^q=R^qf_*\Z.
\]
Proper base change identifies the stalk of $R^q$ at a special point with
the cohomology of the corresponding fibre.  At $p_0$, the map to the
nearby invariant lattice is the specialization map of
\cref{thm:cusp-specialization}.  At $p_j$, the local tube retracts to
$S_j$, and the map to a nearby fibre is the pullback along the covering
$F\to S_j$.  These maps are injective by
\cref{prop:finite-specialization,thm:cusp-specialization}; hence $R^q$ is
the subsheaf of $j_*L^q$ determined by the three displayed stalk lattices.

\begin{lemma}[Cohomology of the extended local systems]
\label{lem:local-system-cohomology}
Let $L$ be a local system of free finitely generated abelian groups on
$B^\circ$, with stalk $M$ and monodromy group $G$.
\begin{enumerate}[label=\textup{(\roman*)}]
\item There is a natural isomorphism
\begin{equation}\label{eq:H2-coinvariants}
 H^2(B,j_*L)\cong M_G.
\end{equation}
It is compatible with pairings of local systems: under
\eqref{eq:H2-coinvariants}, cup product by an invariant element is the
induced action on coinvariants.
\item Let $t_1,t_2$ be the clockwise generators of
$\pi_1(B^\circ)$, let $t_0=(t_1t_2)^{-1}$, and write $T_i$ for their
actions on $M$.  Then
\begin{equation}\label{eq:parabolic-cohomology}
 H^1(B,j_*L)
 \cong
 \frac{(T_1-I)M\cap(T_0-I)M}{(T_1-I)M^{T_2}}.
\end{equation}
\end{enumerate}
\end{lemma}

\begin{proof}
The quotient $j_*L/j_!L$ is supported at the three punctures and has no
cohomology in degree two.  Therefore
\[
 H^2(B,j_*L)\cong H_c^2(B^\circ,L)
 \cong H_0(B^\circ,L)\cong M_G
\]
by Poincaré duality with local coefficients \cite[Ch.~V]{Bredon}.
Naturality of cap product gives compatibility with coefficient pairings.

A class in $H^1(B,j_*L)$ may be represented by an inhomogeneous group
cocycle on $\pi_1(B^\circ)$.  Extension across the puncture $p_i$ means
that its restriction to the local cyclic subgroup is a coboundary; this
is equivalent to the condition $c(t_i)\in(T_i-I)M$.  Thus the relevant
classes are precisely the parabolic cocycles.  Subtracting a coboundary,
arrange $c(t_2)=0$; the remaining coboundaries are those defined by
elements of $M^{T_2}$.  Put
$x=c(t_1)$.  The cocycle identity and $t_0=(t_1t_2)^{-1}$ give
$c(t_0)=-T_0x$.  Hence the parabolic conditions are
$x\in(T_1-I)M\cap(T_0-I)M$, and the residual coboundaries change $x$ by
$(T_1-I)M^{T_2}$.  This proves \eqref{eq:parabolic-cohomology}.
\end{proof}

\begin{proposition}\label{prop:Rq-cohomology}
The quotient $j_*L^q/R^q$ is a skyscraper sheaf with stalks
\[
\begin{array}{c|ccc}
q&p_0&p_1&p_2\\ \hline
1&0&\Z/3&\Z/4\\
2&0&0&\Z/2\\
3&0&0&\Z/2\\
4&0&\Z/3&\Z/4.
\end{array}
\]
Consequently
\begin{align}
 H^0(B,R^1)&=12\Z\gamma,
 &H^0(B,R^2)&=2\Z q,
 \label{eq:H0R12}\\
 H^0(B,R^3)&=2\Z\gamma uw,
 &H^0(B,R^4)&=12\Z\vol.
 \label{eq:H0R34}
\end{align}
Furthermore
\begin{align}
 H^2(B,R^1)&=\Z[\delta],
 &H^2(B,R^2)&=\Z[\gamma\delta],
 \label{eq:H2R12}\\
 H^2(B,R^3)&=\Z[uw\delta],
 &H^2(B,R^4)&=\Z[\vol],
 \label{eq:H2R34}
\end{align}
and
\begin{equation}\label{eq:H1-vanishing}
 H^1(B,R^1)=H^1(B,R^2)=0.
\end{equation}
The identifications in degree two are compatible with cup products: an
invariant class acts on coinvariants by exterior product.
\end{proposition}

\begin{proof}
The skyscraper table is exactly
\cref{prop:finite-specialization,thm:cusp-specialization}.  Taking global
invariants and imposing the local divisibilities gives
\eqref{eq:H0R12}--\eqref{eq:H0R34}, using
\cref{prop:linear-algebra}(vi).

Since the skyscraper quotient $j_*L^q/R^q$ has no degree-two
cohomology, \cref{lem:local-system-cohomology}(i) and the coinvariant
calculations of \cref{prop:linear-algebra}(vi) give
\eqref{eq:H2R12}--\eqref{eq:H2R34} with the stated compatibility.

The maps from the global invariant lattices to the skyscraper quotients
are surjective in degrees one and two: in degree one, $\gamma$ maps to
$(1,1)\in\Z/3\oplus\Z/4$, and in degree two, $q$ maps to the generator of
$\Z/2$.  Hence $H^1(B,R^q)=H^1(B,j_*L^q)$ for $q=1,2$.  Apply
\cref{lem:local-system-cohomology}(ii).  For $M=V$,
\[
 (T_1-I)V=\langle2\gamma+u+w,2\gamma-u\rangle,
 \qquad
 (T_0-I)V=\langle\gamma,u\rangle,
\]
whose intersection is $\Z(2\gamma-u)$.  Since
$(T_1-I)\psi_2=-(2\gamma-u)$, this equals $(T_1-I)V^{T_2}$.  For
$M=\bigwedge^2V$,
\[
 (\textstyle\bigwedge^2T_0-I)M
 =\langle\gamma u,\gamma w+u\delta\rangle,
\]
while $(\bigwedge^2T_1-I)M$ is generated by
\[
 \gamma u,\ \gamma w,\ 2u\delta+w\delta,
 \ uw+u\delta-w\delta-6\gamma\delta.
\]
The intersection is $\Z\gamma u$.  Since $\gamma\psi_2$ is fixed by
$\bigwedge^2T_2$ and
$(\bigwedge^2T_1-I)(\gamma\psi_2)=\gamma u$, this is
$(\bigwedge^2T_1-I)M^{T_2}$.  Formula
\eqref{eq:parabolic-cohomology} therefore gives the required vanishing.
\end{proof}

\subsection{The Leray differential}
Consider the integral Leray spectral sequence
\[
 E_2^{p,q}=H^p(B,R^qf_*\Z)\Longrightarrow H^{p+q}(X;\Z).
\]
Since $B$ has real dimension two, the only possible nonzero differentials
are
\[
 d_2^{0,q}:E_2^{0,q}\longrightarrow E_2^{2,q-1}.
\]

\begin{proposition}\label{prop:leray-differentials}
After choosing the sign of the generator
$\omega\in H^2(B;\Z)$, one has
\begin{align}
 d_2(12\gamma)&=\omega,
 \label{eq:d2-gamma}\\
 d_2(2q)&=\omega[\delta],
 \label{eq:d2-q}\\
 d_2(2\gamma uw)&=\omega[\gamma\delta].
 \label{eq:d2-gamma-uw}
\end{align}
In particular all three maps are isomorphisms of infinite cyclic groups.
\end{proposition}

\begin{proof}
The first differential is the obstruction to lifting the global section
$12\gamma$ of $R^1$ to a class in $H^1(X;\Z)$.  We compute it by local
lifts.  Let $U=D_0\sqcup D_1\sqcup D_2$ be a disjoint union of small
discs and let $V$ be a neighbourhood of the complementary pair of pants,
so that $U\cap V$ is a disjoint union of three annuli.  For each of the three discs and for $V$, the five-term Leray exact
sequence contains
\[
 H^1(f^{-1}(Y);\Z)\longrightarrow H^0(Y,R^1)
 \longrightarrow H^2(Y;\Z),
\]
where $Y$ denotes the corresponding base region.  The last group vanishes.
Thus a section of $R^1$ lifts to a degree-one class on the inverse image
exactly when it lies in the appropriate local specialization lattice.
This holds for $12\gamma$ by
\cref{prop:finite-specialization,thm:cusp-specialization}.
On each annulus the two lifts have the same image in $H^0(R^1)$.  The
five-term exact sequence for the annulus therefore identifies their
difference with the pullback of a unique class in $H^1$ of the annulus.
These three classes form a Čech cocycle, and its image under the
Mayer--Vietoris isomorphism
\[
 H^1(U\cap V;\Z)/\im H^1(V;\Z)
 \xrightarrow{\sim}H^2(B;\Z)
\]
is exactly $d_2(12\gamma)$.  Evaluating a difference on the oriented
annular meridian identifies $H^1(U\cap V;\Z)$ with $\Z^3$; the image of
$H^1(V;\Z)$ is the sum-zero sublattice, so the class of
$(n_0,n_1,n_2)$ is $\pm(n_0+n_1+n_2)\omega$.

Over $D_0$, choose the lift to have value zero on the toric meridian,
which is null-homotopic in $N_0$.  Over $D_j$, if $\Theta_j$ is a lift,
the relation $\wh g_j^{m_j}=t_{v_j}$ gives
\[
 m_j\Theta_j(\wh g_j)=12\gamma(v_j).
\]
Thus the meridian values of the local lifts are
\[
 4\quad\text{at }p_1,
 \qquad
 -3\quad\text{at }p_2,
 \qquad
 0\quad\text{at }p_0.
\]
Let the lift over $V$ have values $a_1,a_2$ on the first two meridians.
The product relation for the three clockwise boundary meridians gives
value $-a_1-a_2$ at the cusp.  The three differences between the disc
lifts and the $V$-lift are therefore
\[
 4-a_1,
 \qquad
 -3-a_2,
 \qquad
 a_1+a_2.
\]
Their sum is $1$.  By the preceding Mayer--Vietoris identification,
choosing the sign of $\omega$ accordingly proves \eqref{eq:d2-gamma}.

Write
\[
 d_2(2q)=a\,\omega[\delta],
 \qquad
 d_2(2\gamma uw)=b\,\omega[\gamma\delta]
\]
for integers $a,b$.  Multiplicativity and the relation
\[
 (12\gamma)(2q)=12(2\gamma uw)
\]
give, using the graded Leibniz rule with the sign $(-1)^1=-1$ and the
identity $[q]=12[\gamma\delta]$ in the degree-two coinvariants,
\[
 12b=24-12a,
 \qquad\text{hence}\qquad b=2-a.
\]
Next $(2q)(2\gamma uw)=0$ because $H^5(F;\Z)=0$.  Under the
coinvariant identifications of \cref{lem:local-system-cohomology}, the two
coefficient pairings
\[
 H^2(B,R^1)\otimes H^0(B,R^3)\longrightarrow H^2(B,R^4),
 \qquad
 H^0(B,R^2)\otimes H^2(B,R^2)\longrightarrow H^2(B,R^4)
\]
have absolute value $2$ on the displayed generators: on representatives,
\[
 \delta\wedge(2\gamma uw)=-2\vol,
 \qquad
 (2q)\wedge(\gamma\delta)=2\vol.
\]
The graded Leibniz rule applied to the zero product therefore gives
$|a|=|b|$.  Since $b=2-a$, one has $a^2=(2-a)^2$, hence $a=b=1$.
This proves
\eqref{eq:d2-q}--\eqref{eq:d2-gamma-uw}.
\end{proof}

\begin{theorem}\label{thm:integral-cohomology}
The integral homology of $X$ is that of $S^6$.
\end{theorem}

\begin{proof}
By \cref{prop:Rq-cohomology,prop:leray-differentials}, the spectral
sequence gives
\[
 H^1(X;\Z)=H^2(X;\Z)=H^3(X;\Z)=0.
\]
Indeed, in total degree one the only term is the kernel of
\eqref{eq:d2-gamma}; in degree two the three graded pieces are the
cokernel of \eqref{eq:d2-gamma}, $H^1(B,R^1)$, and the kernel of
\eqref{eq:d2-q}; in degree three they are the cokernel of
\eqref{eq:d2-q}, $H^1(B,R^2)$, and the kernel of
\eqref{eq:d2-gamma-uw}.

Since $\pi_1(X)=1$ by \cref{thm:pi1}, one has $H_1(X;\Z)=0$.
The universal coefficient theorem and $H^2(X;\Z)=0$ imply
$\Hom(H_2(X),\Z)=0$, so the finitely generated group $H_2(X)$ is torsion.
The group $\operatorname{Ext}(H_2(X),\Z)$ injects into $H^3(X;\Z)=0$;
for a finite abelian group it is isomorphic to the group itself.  Hence
$H_2(X;\Z)=0$.  Integral Poincaré duality now gives
\[
 H_3(X;\Z)=H_4(X;\Z)=H_5(X;\Z)=0.
\]
Finally $H_0(X;\Z)=H_6(X;\Z)=\Z$ because $X$ is connected and oriented.
\end{proof}


\section{Recognition of the smooth manifold}

\begin{lemma}\label{lem:homology-sphere}
Let $M$ be a closed simply connected $n$-manifold, $n\ge2$, with the
integral homology of $S^n$.  Then $M$ is homotopy equivalent to $S^n$.
\end{lemma}

\begin{proof}
If $k<n$ were the least integer with $\pi_k(M)\ne0$, the Hurewicz
theorem would give $\pi_k(M)\cong H_k(M;\Z)=0$, a contradiction.
Thus $\pi_i(M)=0$ for $1\le i<n$, and the Hurewicz theorem in degree
$n$ gives an isomorphism $\pi_n(M)\to H_n(M;\Z)\cong\Z$.  Choose
$u:S^n\to M$ representing a class mapping to a generator.  The map $u$
is an isomorphism on integral homology.  Since both spaces have the
homotopy type of CW complexes and are simply connected, the homology
Whitehead theorem implies that $u$ is a homotopy equivalence
\cite[Ch.~4]{Hatcher}.
\end{proof}

\begin{proof}[Proof of \cref{thm:main}]
By \cref{thm:pi1,thm:integral-cohomology,lem:homology-sphere}, the closed
smooth manifold underlying $X$ is a homotopy six-sphere.  For $n\ge5$, the group $\Theta_n$ consists of the oriented
$h$-cobordism classes of homotopy $n$-spheres.  In dimensions at least five, the smooth
$h$-cobordism theorem makes two such spheres diffeomorphic whenever they
represent the same class \cite{Milnor}.  Kervaire and Milnor computed
$\Theta_6=0$ \cite{KM}.  Hence $X$ is
orientation-preservingly diffeomorphic to the standard $S^6$.  Transporting
the complex atlas of $X$ through such a diffeomorphism gives an integrable
complex structure on $S^6$.
\end{proof}



\appendix
\section{Integral nearby cycles at the cusp}\label{sec:nearby-cycles}

Let $\psi=R\psi_{f_0}\Z$ be the nearby-cycle complex on $W$ in the
analytic topology \cite[Ch.~4]{Dimca}.  Proper base change gives
\[
 \mathbb H^*(W,\psi)\cong H^*(F;\Z),
\]
and the edge map
$H^q(W;\Z)\to H^q(F;\Z)$ is the specialization map.

\subsection{The cohomology sheaves of nearby cycles}
At a point of $W$, the map $f_0$ has one of the local forms
\[
 z_0,
 \qquad
 z_0z_1,
 \qquad
 z_0z_1z_2.
\]
Writing $z_i=r_ie^{i\theta_i}$, the radial constraint in the Milnor
fibre of $z_0\cdots z_k=t$ is an open simplex, while the phase constraint
$\theta_0+\cdots+\theta_k=\arg t$ is a $k$-torus.  Thus the three local
Milnor fibres are homotopy equivalent to a point, $S^1$, and $T^2$.
Therefore
\[
 \cH^0\psi=\Z_W,
 \qquad
 \cH^1\psi=\cV,
 \qquad
 \cH^2\psi=\Z_P\oplus\Z_Q,
\]
where $P,Q$ are the two triple points and $\cV$ is supported on the double
locus $D=D_1\cup D_2\cup D_3$.  Write $\iota_k:D_k\hookrightarrow W$ for
the inclusions.

Number the six sides of the hexagon cyclically by
$C_1,\ldots,C_6$ and let $D_k$ be the common image of $C_k$ and
$C_{k+3}$, for $k=1,2,3$.  The normalization distinguishes the two
branches globally along $D_k\setminus\{P,Q\}$, so the vanishing-cycle
local system there is trivial.  Orient its generator by increasing
argument on the branch coming from $C_k$.  At either triple
point the three restriction values
$n_1,n_2,n_3$ satisfy
\[
 n_1-n_2+n_3=0.
\]
Indeed, in the local model $z_0z_1z_2=t$, the three vanishing circles have
homology classes $e_0-e_1$, $e_0-e_2$, $e_1-e_2$ in
$\{(a_0,a_1,a_2)\in\Z^3:\sum a_i=0\}$.  The global signs are as follows.
The three preimages of $P$ are
$C_1\cap C_2$, $C_3\cap C_4$, and $C_5\cap C_6$; call the corresponding
normalization branches $b_0,b_1,b_2$.  Thus the ordered branch pairs
defining $D_1,D_2,D_3$ are
$(b_0,b_1),(b_0,b_2),(b_1,b_2)$, giving the three differences above.
The preimages of $Q$ are
$C_2\cap C_3$, $C_4\cap C_5$, and $C_6\cap C_1$; call the corresponding
branches $c_0,c_1,c_2$.  The ordered pairs are then
$(c_2,c_1),(c_0,c_1),(c_0,c_2)$.  The corresponding differences are
$e_2-e_1,e_0-e_1,e_0-e_2$, and again satisfy
$n_1-n_2+n_3=0$.  Hence there is an exact sequence
\begin{equation}\label{eq:V-resolution}
 0\longrightarrow\cV\longrightarrow
 \bigoplus_{k=1}^3\iota_{k*}\Z_{D_k}
 \xrightarrow{d}
 \Z_P\oplus\Z_Q\longrightarrow0,
 \qquad
 d(n_1,n_2,n_3)=(n_1-n_2+n_3,n_1-n_2+n_3).
\end{equation}

Exactness is stalkwise.  Away from $P,Q$ the restriction map is an
isomorphism onto the single local branch of the double locus.  At a triple
point the stalk of $\cV$ is
$\Hom(\{(a_0,a_1,a_2)\in\Z^3:\sum a_i=0\},\Z)$, and the three
restriction homomorphisms are evaluation on the three differences above.
Their image is exactly the kernel of $n_1-n_2+n_3$, and that linear form
is primitive, so the cokernel is $\Z$.

\subsection{The cohomology of the central fibre}
Let $\nu:dP_6\to W$ be the normalization.  For a locally constant
function on the normalization, define its $k$-th difference to be the
restriction to $C_k$ minus the restriction to $C_{k+3}$, transported to
$D_k$.  With this difference map and the map $d$ of
\eqref{eq:V-resolution}, the normalization complex is
\begin{equation}\label{eq:W-normalization-resolution}
 0\longrightarrow\Z_W\longrightarrow\nu_*\Z_{dP_6}
 \longrightarrow\bigoplus_{k=1}^3\iota_{k*}\Z_{D_k}
 \xrightarrow{d}\Z_P\oplus\Z_Q\longrightarrow0.
\end{equation}

This is again exact on stalks.  At a smooth point it is
$0\to\Z\to\Z\to0$; at a general double point it is the augmented
cochain complex of an interval,
$0\to\Z\to\Z^2\to\Z\to0$; and at a triple point it is the
augmented cochain complex of a two-simplex,
$0\to\Z\to\Z^3\to\Z^3\to\Z\to0$.
The signs are the same as in \eqref{eq:V-resolution} because both
complexes use the same ordering of the local branches.  Since $\nu$ is
finite, proper base change gives $R^r\nu_*\Z=0$ for $r>0$; the same
vanishing is immediate for the closed embeddings $\iota_k$ and for the
point inclusions.  Thus the sheaf cohomology of the terms in
\eqref{eq:W-normalization-resolution} is the ordinary cohomology of
$dP_6$, the curves $D_k$, and the points $P,Q$.

The map in degree two
$H^2(dP_6;\Z)\to\bigoplus_kH^2(D_k;\Z)=\Z^3$ sends a divisor to the
differences of its intersection numbers with opposite sides of the
hexagon.  Writing $dP_6$ as the blow-up of $\PP^2$ at three points, with
basis $H,E_1,E_2,E_3$, the six sides are
\[
 E_1,
 H-E_1-E_2,
 E_2,
 H-E_2-E_3,
 E_3,
 H-E_1-E_3.
\]
For $L=aH-\sum b_iE_i$, the difference map is
\[
 L\longmapsto(s,-s,s),
 \qquad s=b_1+b_2+b_3-a.
\]
Its image is the saturated subgroup $\Z(1,-1,1)$.
The hypercohomology spectral sequence is explicit.  In sheaf-cohomology degree zero,
the complex of global sections is
\[
 \Z\longrightarrow\Z^3\xrightarrow{(a_1,a_2,a_3)\mapsto
 (a_1-a_2+a_3)(1,1)}\Z^2,
\]
with first map zero.  Its cohomology has ranks $1,2,1$ in complex degrees
$0,1,2$; this is the cohomology of the dual cell complex, a torus with one
vertex, three edges and two faces.  The degree-two row is
\[
 H^2(dP_6;\Z)=\Z^4\longrightarrow\Z^3,
\]
with saturated rank-one image, hence kernel $\Z^3$ and cokernel
$\Z^2$.  The degree-four row is $H^4(dP_6;\Z)=\Z$.  The resolution has only
three nonzero terms after $\Z_W$, in complex degrees $0,1,2$.  Hence
$d_r=0$ for $r\geq3$, while the only possible $d_2$ has target in an odd
cohomology row and therefore vanishes.  The spectral sequence degenerates
at $E_2$.  All surviving groups are free, so there is no extension
torsion.  Therefore
\begin{equation}\label{eq:W-cohomology-ranks}
 H^q(W;\Z)=\Z,\ \Z^2,\ \Z^4,\ \Z^2,\ \Z
 \quad(q=0,1,2,3,4),
\end{equation}
with no torsion.

From \eqref{eq:V-resolution} and $H^*(\PP^1;\Z)=(\Z,0,\Z)$ one obtains
\begin{equation}\label{eq:V-cohomology}
 H^q(W;\cV)=\Z^2,\ \Z,\ \Z^3,\ 0
 \quad(q=0,1,2,3).
\end{equation}

\subsection{The nearby-cycle spectral sequence}
The spectral sequence
\[
 E_2^{a,b}=H^a(W;\cH^b\psi)\Longrightarrow H^{a+b}(F;\Z)
\]
has the following free $E_2$-terms, displayed by rank:
\[
\begin{array}{c|ccccc}
b=2&2&0&0&0&0\\
b=1&2&1&3&0&0\\
b=0&1&2&4&2&1\\ \hline
&a=0&a=1&a=2&a=3&a=4.
\end{array}
\]
Let
\begin{align*}
 \alpha&=\rk(d_2:E_2^{0,1}\to E_2^{2,0}),
 &\beta&=\rk(d_2:E_2^{1,1}\to E_2^{3,0}),\\
 \gamma&=\rk(d_2:E_2^{2,1}\to E_2^{4,0}),
 &\delta&=\rk(d_2:E_2^{0,2}\to E_2^{2,1}),\\
 \epsilon&=\rk(d_3:E_3^{0,2}\to E_3^{3,0}).
\end{align*}
These are the only possible nonzero differentials.  Since the Betti
numbers of $F=T^4$ are $(1,4,6,4,1)$, total-degree bookkeeping gives
\[
 4=2+(2-\alpha),
 \qquad
 1=1-\gamma,
\]
so $\alpha=\gamma=0$.  In total degrees two and three one obtains
\[
 6=(4-\alpha)+(1-\beta)+(2-\delta-\epsilon),
\]
\[
 4=(2-\beta-\epsilon)+(3-\gamma-\delta).
\]
Hence
\[
 \beta+\delta+\epsilon=1.
\]
Thus specialization is already injective in degrees $0,1,2,4$, and
exactly one rank-one differential remains among
\[
 E_2^{1,1}\to E_2^{3,0},
 \qquad
 E_2^{0,2}\to E_2^{2,1},
 \qquad
 E_3^{0,2}\to E_3^{3,0}.
\]
It remains to control specialization in degree three.

\begin{lemma}[Specialization with supports]\label{lem:specialization-supports}
Let $K$ be a compact subset of the smooth locus of $W$.  After shrinking
the base disc, there are a neighbourhood $K'$ of $K$ in $W$ and an open
embedding
\[
 \Theta:K'\times\Delta_r\longrightarrow N_0,
 \qquad f_0(\Theta(w,t))=t,
 \qquad \Theta(w,0)=w.
\]
If a compact group acts on $N_0$, preserves $f_0$ and $K$, then $K'$ and
$\Theta$ may be chosen equivariantly.  For
$K_t=\Theta(K\times\{t\})$, the composite
\[
 H^q(W,W\setminus K)\longrightarrow H^q(W)
 \xrightarrow{\operatorname{sp}^q}H^q(F_t)
\]
equals, under transport by $\Theta$,
\[
 H^q(F_t,F_t\setminus K_t)\longrightarrow H^q(F_t).
\]
\end{lemma}

\begin{proof}
Choose a relatively compact neighbourhood of $K$ on which $f_0$ is a
submersion and a Riemannian metric there.  If a compact group is present,
replace the neighbourhood by its group saturation and average the metric
over the group.  Let $\xi_1,\xi_2$ be the horizontal lifts of
$\partial/\partial\Re t$ and $\partial/\partial\Im t$.  For
$t\in\C$ set
\[
 \xi_t=(\Re t)\xi_1+(\Im t)\xi_2.
\]
Then $df_0(\xi_t)=t$.  For $r>0$ sufficiently small, the time-one flow of
$\xi_t$ is defined on a fixed neighbourhood $K'$ of $K$ for all
$|t|<r$.  Put
\[
 \Theta(w,t)=\operatorname{Fl}^{1}_{\xi_t}(w).
\]
Along the flow, $f_0$ has derivative $t$, so
$f_0(\Theta(w,t))=t$.  At $t=0$ the map is the identity.  If
$\Theta(w,t)=\Theta(w',t')$, applying $f_0$ gives $t=t'$, and injectivity
of the time-one flow of the fixed vector field $\xi_t$ then gives
$w=w'$.  The differential of $\Theta$ is invertible along
$K\times\{0\}$; after shrinking $K'$ and $r$, it is therefore an open
embedding.  The construction is equivariant when the metric is.

Let $Z=\Theta(K\times\Delta_r)$.  The restriction
$\Theta|_{K\times\Delta_r}$ is proper: for a compact subset $L$ of
$f_0^{-1}(\Delta_r)$, its inverse image is a closed subset of the compact
set $K\times f_0(L)$.  Hence $Z$ is closed in $f_0^{-1}(\Delta_r)$.  Excision identifies the relative
cohomology of $(f_0^{-1}(\Delta_r),f_0^{-1}(\Delta_r)\setminus Z)$ with
that of the product pair
$(K'\times\Delta_r,(K'\setminus K)\times\Delta_r)$.  Restriction to the
slices $t=0$ and $t\ne0$ is an isomorphism by homotopy invariance.  The
specialization map is obtained by restricting from a sufficiently small
tube to a nearby fibre, so these identifications give the asserted
commutative diagram.
\end{proof}

Let $K\cong T^2$ be an orbit of the compact torus
$T_f=\Lambda_{\mathrm{tor}}\otimes S^1$ in the open torus orbit of $W$.
The open orbit is $(\C^*)^2$, so the normal bundle of $K$ in the smooth
four-manifold $W^\circ$ is trivial.  Apply the equivariant part of
\cref{lem:specialization-supports}; then $K_t$ is a $T_f$-orbit in the
nearby torus.  The composite
\[
 H^3(W,W\setminus K)\longrightarrow H^3(W)
 \xrightarrow{\operatorname{sp}^3}H^3(F)
\]
is the Gysin map of the transported orbit $K_t\subset F$.  The Thom
isomorphism identifies the source with $H^1(K)$.

Under the exponential coordinates of
\cref{prop:cusp-identification}, the $T_f$-action on a nearby fibre is
translation by
$(\Lambda_{\mathrm{tor}}\otimes\R)/\Lambda_{\mathrm{tor}}$.
Consequently $K_t$ corresponds to the saturated sublattice
$\Lambda_{\mathrm{tor}}\subset\Lambda$.  Choose a complementary
sublattice $\Lambda'$, so
$F\cong K_t\times T_B$ with
$T_B=(\Lambda'\otimes\R)/\Lambda'$.  Under this product decomposition,
the Gysin map is
\[
 H^1(K_t)\longrightarrow H^3(F),
 \qquad a\longmapsto a\smile[T_B]^*.
\]
Its image is a direct summand of the Künneth decomposition and is therefore
a saturated rank-two subgroup.  The composite through $H^3(W)$ has rank
two, while $H^3(W)$ is free of rank two; hence
$\operatorname{sp}^3$ is injective.  Its image contains the saturated
rank-two Gysin image and has the same rank, so the two images coincide.
In particular the specialization image is saturated.  Hence both
possible differentials landing in $E^{3,0}$ vanish:
$\beta=\epsilon=0$.  Therefore $\delta=1$, and the unique rank-one
differential is $E_2^{0,2}\to E_2^{2,1}$.

It follows that specialization is injective in every degree.  Its image is
saturated.  In degree one,
$H^1(F)/\im(\operatorname{sp}^1)=E_\infty^{0,1}=\Z^2$ is free.  In
degree two, the filtration quotient
$H^2(F)/\im(\operatorname{sp}^2)$ has successive graded pieces
$E_\infty^{1,1}=\Z$ and
$E_\infty^{0,2}=\ker(d_2:\Z^2\to\Z^3)$, both free; hence the whole
quotient is free.  Degree three was just proved, and in degree four the
edge map is an isomorphism.  Every specialized class is $T_0$-invariant.
By
\cref{prop:linear-algebra}(vi), the invariant ranks are
$1,2,4,2,1$, exactly the ranks in \eqref{eq:W-cohomology-ranks}.  Since the
invariant lattices are saturated, the saturated specialization image of
the same rank equals the full invariant lattice.  This proves
\cref{thm:cusp-specialization}.


\section{Integral matrix calculations}\label{sec:matrix-tables}
All groups are written in the ordered bases introduced in
\cref{prop:linear-algebra}.

The images on the lattice are
\[
 (A_1-I)\Lambda
 =\langle-\wh u-\wh w+\wh\delta,
             \wh u-2\wh w\rangle,
\]
\[
 (A_2-I)\Lambda
 =\langle-\wh u+\wh w,
             -\wh u-\wh w+\wh\delta\rangle.
\]
The relevant $2\times2$ minors have greatest common divisor one, so both
subgroups are saturated.  On $V$,
\[
 (T_1-I)V=\langle2\gamma+u+w,2\gamma-u\rangle,
 \qquad
 (T_2-I)V=\langle-3\gamma+u,u+w\rangle.
\]
Their sum is the saturated lattice $\langle\gamma,u,w\rangle$.

In the degree-two basis
$(\gamma u,\gamma w,\gamma\delta,uw,u\delta,w\delta)$, direct
application of the exterior-square matrices gives
\[
 (\textstyle\bigwedge^2T_1-I)\bigwedge^2V
 =\langle
 \gamma u,\gamma w,
 2u\delta+w\delta,
 uw+u\delta-w\delta-6\gamma\delta
 \rangle,
\]
\[
 (\textstyle\bigwedge^2T_0-I)\bigwedge^2V
 =\langle\gamma u,\gamma w+u\delta\rangle.
\]
Together, the images of $\bigwedge^2T_1-I$ and
$\bigwedge^2T_2-I$ generate
\[
 \langle\gamma u,\gamma w,u\delta,w\delta,
 uw-6\gamma\delta\rangle.
\]
Adjoining $\gamma\delta$ to these five generators gives a basis of
$\bigwedge^2V$; hence this rank-five lattice is saturated.  This proves
$(\bigwedge^2V)_G=\Z[\gamma\delta]$ and
$[uw]=6[\gamma\delta]$.  Solving the kernel equations gives
\[
 (\textstyle\bigwedge^2V)^{T_0}
 =\langle\gamma u,\gamma\delta,uw,\gamma w-u\delta\rangle,
\]
\[
 (\textstyle\bigwedge^2V)^{T_j}
 =\langle\gamma\psi_j,q\rangle.
\]
The intersection of the two lattices for $j=1,2$ is $\Z q$: comparison of
the $\gamma w$- and $\gamma u$-coefficients forces the coefficient of
$\gamma\psi_j$ to vanish.  Thus
$(\bigwedge^2V)^G=\Z q$.  The products are
\[
 (\gamma\psi_1)q=3\vol,
 \qquad
 (\gamma\psi_2)q=2\vol,
 \qquad
 q^2=12\vol.
\]

The invariant generator $\vol\in\bigwedge^4V$ identifies
$\bigwedge^3V$ equivariantly with the contragredient lattice $\Lambda$.  Hence
$(\bigwedge^3V)^G=\Z\gamma uw$ and
$(\bigwedge^3V)_G=\Z[uw\delta]$, with no torsion, by
\cref{prop:linear-algebra}(v).  Directly,
\[
 (\textstyle\bigwedge^3V)^{T_0}
 =\langle\gamma uw,\gamma u\delta\rangle,
\]
\[
 (\textstyle\bigwedge^3V)^{T_1}
 =\langle\gamma uw,-uw\delta+2\gamma w\delta+4\gamma u\delta\rangle,
\]
\[
 (\textstyle\bigwedge^3V)^{T_2}
 =\langle\gamma uw,-uw\delta+3\gamma w\delta+3\gamma u\delta\rangle.
\]
Pairing these bases with $(\gamma,\psi_1)$ and
$(\gamma,\psi_2)$ gives determinants $-3$ and $-2$, respectively.



\begin{thebibliography}{99}

\bibitem{AnthropicProof}
An AI model internal to Anthropic,
\emph{A compact complex threefold fibred by tori over the projective line, and the six-sphere},
manuscript communicated by Levent Alpoge, 2026,
available at \url{https://alpo.ge/s6.pdf}.

\bibitem{Beardon}
A.~F.~Beardon,
\emph{The Geometry of Discrete Groups},
Graduate Texts in Mathematics 91, Springer, 1983.

\bibitem{Bredon}
G.~E.~Bredon,
\emph{Sheaf Theory}, 2nd ed.,
Graduate Texts in Mathematics 170, Springer, 1997.

\bibitem{Dimca}
A.~Dimca,
\emph{Sheaves in Topology},
Universitext, Springer, 2004.

\bibitem{Fulton}
W.~Fulton,
\emph{Introduction to Toric Varieties},
Annals of Mathematics Studies 131, Princeton University Press, 1993.

\bibitem{Hatcher}
A.~Hatcher,
\emph{Algebraic Topology},
Cambridge University Press, 2002.

\bibitem{KM}
M.~Kervaire and J.~Milnor,
\emph{Groups of homotopy spheres I},
Ann.\ of Math. 77 (1963), 504--537.

\bibitem{Milnor}
J.~Milnor,
\emph{Lectures on the $h$-cobordism theorem},
Princeton University Press, 1965.

\bibitem{Serre}
J.-P.~Serre,
\emph{A Course in Arithmetic},
Graduate Texts in Mathematics 7, Springer, 1973.

\end{thebibliography}



\end{document}
